% Options for packages loaded elsewhere
% Options for packages loaded elsewhere
\PassOptionsToPackage{unicode}{hyperref}
\PassOptionsToPackage{hyphens}{url}
\PassOptionsToPackage{dvipsnames,svgnames,x11names}{xcolor}
%
\documentclass[
  german,
  10pt,
  twoside]{article}
\usepackage{xcolor}
\usepackage[top=2.5cm,bottom=2.5cm,left=2cm,right=2cm]{geometry}
\usepackage{amsmath,amssymb}
\setcounter{secnumdepth}{5}
\usepackage{iftex}
\ifPDFTeX
  \usepackage[T1]{fontenc}
  \usepackage[utf8]{inputenc}
  \usepackage{textcomp} % provide euro and other symbols
\else % if luatex or xetex
  \usepackage{unicode-math} % this also loads fontspec
  \defaultfontfeatures{Scale=MatchLowercase}
  \defaultfontfeatures[\rmfamily]{Ligatures=TeX,Scale=1}
\fi
\usepackage{lmodern}
\ifPDFTeX\else
  % xetex/luatex font selection
\fi
% Use upquote if available, for straight quotes in verbatim environments
\IfFileExists{upquote.sty}{\usepackage{upquote}}{}
\IfFileExists{microtype.sty}{% use microtype if available
  \usepackage[]{microtype}
  \UseMicrotypeSet[protrusion]{basicmath} % disable protrusion for tt fonts
}{}
\usepackage{setspace}
\makeatletter
\@ifundefined{KOMAClassName}{% if non-KOMA class
  \IfFileExists{parskip.sty}{%
    \usepackage{parskip}
  }{% else
    \setlength{\parindent}{0pt}
    \setlength{\parskip}{6pt plus 2pt minus 1pt}}
}{% if KOMA class
  \KOMAoptions{parskip=half}}
\makeatother
% Make \paragraph and \subparagraph free-standing
\makeatletter
\ifx\paragraph\undefined\else
  \let\oldparagraph\paragraph
  \renewcommand{\paragraph}{
    \@ifstar
      \xxxParagraphStar
      \xxxParagraphNoStar
  }
  \newcommand{\xxxParagraphStar}[1]{\oldparagraph*{#1}\mbox{}}
  \newcommand{\xxxParagraphNoStar}[1]{\oldparagraph{#1}\mbox{}}
\fi
\ifx\subparagraph\undefined\else
  \let\oldsubparagraph\subparagraph
  \renewcommand{\subparagraph}{
    \@ifstar
      \xxxSubParagraphStar
      \xxxSubParagraphNoStar
  }
  \newcommand{\xxxSubParagraphStar}[1]{\oldsubparagraph*{#1}\mbox{}}
  \newcommand{\xxxSubParagraphNoStar}[1]{\oldsubparagraph{#1}\mbox{}}
\fi
\makeatother


\usepackage{graphicx}
\makeatletter
\newsavebox\pandoc@box
\newcommand*\pandocbounded[1]{% scales image to fit in text height/width
  \sbox\pandoc@box{#1}%
  \Gscale@div\@tempa{\textheight}{\dimexpr\ht\pandoc@box+\dp\pandoc@box\relax}%
  \Gscale@div\@tempb{\linewidth}{\wd\pandoc@box}%
  \ifdim\@tempb\p@<\@tempa\p@\let\@tempa\@tempb\fi% select the smaller of both
  \ifdim\@tempa\p@<\p@\scalebox{\@tempa}{\usebox\pandoc@box}%
  \else\usebox{\pandoc@box}%
  \fi%
}
% Set default figure placement to htbp
\def\fps@figure{htbp}
\makeatother



\ifLuaTeX
\usepackage[bidi=basic]{babel}
\else
\usepackage[bidi=default]{babel}
\fi
% get rid of language-specific shorthands (see #6817):
\let\LanguageShortHands\languageshorthands
\def\languageshorthands#1{}
\ifLuaTeX
  \usepackage[german]{selnolig} % disable illegal ligatures
\fi


\setlength{\emergencystretch}{3em} % prevent overfull lines

\providecommand{\tightlist}{%
  \setlength{\itemsep}{0pt}\setlength{\parskip}{0pt}}



 
\usepackage[style=authoryear,backend=biber,style=authoryear-comp,uniquename=false,uniquelist=false]{biblatex}
\addbibresource{references.bib}


\usepackage{booktabs}
\usepackage{longtable}
\usepackage{array}
\usepackage{multirow}
\usepackage{wrapfig}
\usepackage{float}
\usepackage{colortbl}
\usepackage{pdflscape}
\usepackage{tabu}
\usepackage{threeparttable}
\usepackage{threeparttablex}
\usepackage[normalem]{ulem}
\usepackage{makecell}
\usepackage{xcolor}
% Font configuration (Statis Sans alternative - use sans-serif throughout)
\renewcommand{\familydefault}{\sfdefault}
\usepackage{helvet}  % Use Helvetica-like font

\usepackage{csquotes}
\usepackage{booktabs}
\usepackage{array}
\usepackage{multirow}
\usepackage{wrapfig}
\usepackage{float}
\usepackage{colortbl}
\usepackage{pdflscape}
\usepackage{threeparttable}
\usepackage{threeparttablex}
\usepackage[normalem]{ulem}
\usepackage{makecell}
\usepackage{xcolor}
\usepackage{multicol}
\usepackage{fancyhdr}
\usepackage{titlesec}
\usepackage{etoolbox}
\usepackage{tabularx}
\usepackage{afterpage}
\usepackage{tikz}
\usepackage{longtable}
\usepackage{totcount}
\usepackage{xstring}
\usepackage{pgfmath}  % For calculations

% Time display with Berlin timezone (+1 hour offset)
\newcount\berlinhr
\newcount\berlinmin
\newcommand{\berlintime}{%
  \berlinhr=\time\relax
  \divide\berlinhr by 60\relax
  \berlinmin=\time\relax
  \multiply\berlinhr by -60\relax
  \advance\berlinmin by \berlinhr\relax
  \berlinhr=\time\relax
  \divide\berlinhr by 60\relax
  \advance\berlinhr by 1\relax
  \ifnum\berlinhr>23\relax
    \advance\berlinhr by -24\relax
  \fi
  \ifnum\berlinhr<10 0\fi\the\berlinhr:\ifnum\berlinmin<10 0\fi\the\berlinmin%
}

% Register counters for total counting
\regtotcounter{table}
\regtotcounter{figure}

% Suppress default maketitle
\renewcommand{\maketitle}{}

% ===== DRAFT MODE CONTROL =====
% Set \draftmodetrue to enable draft mode (metadata page + ENTWURF header)
% Set \draftmodefalse to disable draft mode (for final version)
\newif\ifdraftmode
\draftmodetrue  % <-- Change to \draftmodefalse to disable draft mode

% Define default commands for counts (will be overridden by filter)
\providecommand{\totalchars}{0}
\providecommand{\totalwords}{0}
\providecommand{\abstractchars}{0}
\providecommand{\abstractwords}{0}
\providecommand{\zusammenfassungchars}{0}
\providecommand{\zusammenfassungwords}{0}
\providecommand{\footnotecountnum}{0}
\providecommand{\citationcountnum}{0}
\providecommand{\schlusselwortercount}{0}
\providecommand{\keywordscount}{0}
\providecommand{\authoronechars}{0}
\providecommand{\authortwochars}{0}

% Define metadata commands from Pandoc variables
\newcommand{\doctitle}{$title$}
\newcommand{\docauthors}{$author-short$}
\newcommand{\docpublication}{$publication$}
\newcommand{\docpublicationissue}{$publication-issue$}
\newcommand{\docshorttitle}{$short-title$}

\ifdraftmode
% Helper function - show tick or cross based on requirements
\usepackage{amssymb}  % for \checkmark
\usepackage{pifont}   % for \ding{55} (cross mark)

% Check if requirement is met (tick) or not (cross)
% Only checks maximum (no minimum) - requirement is an upper limit
\newcommand{\checkreq}[2]{%
  % #1 = actual value (formatted with spaces like "8 452")
  % #2 = required value (maximum allowed)
  % Remove spaces and compare
  \StrDel{#1}{ }[\cleanvalue]%
  \IfStrEq{\cleanvalue}{0}{%
    % If zero, show cross
    \textcolor{red}{\ding{55}}%
  }{%
    % Check if value exceeds maximum (with 20% tolerance)
    \newcount\actualval%
    \newcount\maxval%
    \actualval=\cleanvalue\relax%
    \maxval=#2\relax%
    \multiply\maxval by 6\relax%
    \divide\maxval by 5\relax%  % 120% = 6/5
    \ifnum\actualval>\maxval\relax
      \textcolor{red}{\ding{55}}%
    \else
      \textcolor{green!60!black}{\checkmark}%
    \fi
  }%
}

% Helper for keyword count check
\newcommand{\checkkeywords}[1]{%
  % #1 = count
  \ifnum#1>5\relax
    \textcolor{red}{\ding{55}}%
  \else%
    \textcolor{green!60!black}{\checkmark}%
  \fi%
}

% Define metadata page command
\newcommand{\metadatapage}{%
  \clearpage
  \thispagestyle{empty}
  \onecolumn
  \vspace*{1.5cm}
  
  % Header section
  {\raggedright
  {\LARGE\bfseries\color{wistaBlue} Dokument-Metadaten}\\[0.5cm]
  
  \begin{tabular}{@{}ll}
    \textbf{Titel:} & \doctitle \\[0.3em]
    \textbf{Autoren:} & \docauthors \\[0.3em]
    \textbf{Veröffentlichung:} & \docpublication\ \docpublicationissue \\[0.3em]
    \textbf{Erstellungsdatum:} & \today\ um \berlintime\ Uhr \\[0.3em]
  \end{tabular}
  
  \vspace{0.5cm}
  {\color{wistaBlue}\rule{\textwidth}{1pt}}
  \par}
  
  \vspace{1cm}
  
  % Dokumentstatistik
  {\large\bfseries\color{wistaDarkGray} Dokumentstatistik}\\[0.5cm]
  
  \small
  \begin{tabular}{@{}lrrrrrl@{}}
    \toprule
    \textbf{Abschnitt} & \textbf{Wörter} & \textbf{Zeichen} & \textbf{Fußn.} & \textbf{Quellen} & \textbf{Vorgabe} & \textbf{Status} \\
    \midrule
    Insgesamt & \textcolor{wistaBlue}{\textbf{\totalwords}} & \textcolor{wistaBlue}{\textbf{\totalchars}} & \textcolor{wistaBlue}{\textbf{\footnotecountnum}} & \textcolor{wistaBlue}{\textbf{\citationcountnum}} & 30\,000 Z. & \checkreq{\totalchars}{30000} \\
    Zusammenfassung & \textcolor{wistaBlue}{\textbf{\zusammenfassungwords}} & \textcolor{wistaBlue}{\textbf{\zusammenfassungchars}} & --- & --- & 800 Z. & \checkreq{\zusammenfassungchars}{800} \\
    Abstract & \textcolor{wistaBlue}{\textbf{\abstractwords}} & \textcolor{wistaBlue}{\textbf{\abstractchars}} & --- & --- & 700 Z. & \checkreq{\abstractchars}{700} \\
    Autor 1 (O. Hauke) & --- & \textcolor{wistaBlue}{\textbf{\authoronechars}} & --- & --- & 350 Z. & \checkreq{\authoronechars}{350} \\
    Autor 2 (A. Kristiansen) & --- & \textcolor{wistaBlue}{\textbf{\authortwochars}} & --- & --- & 350 Z. & \checkreq{\authortwochars}{350} \\
    \midrule
    \sectionrows
    \bottomrule
  \end{tabular}
  \normalsize
  
  \vspace{0.5cm}
  
  % Keywords list with count validation
  {\large\bfseries\color{wistaDarkGray} Schlüsselwörter ($\leq$ 5) \checkkeywords{\schlusselwortercount}}\\[0.3cm]
  \small
  \schlusselworterlist
  \normalsize
  
  \vspace{0.5cm}
  
  {\large\bfseries\color{wistaDarkGray} Keywords ($\leq$ 5) \checkkeywords{\keywordscount}}\\[0.3cm]
  \small
  \keywordslist
  \normalsize
  
  \vspace{1cm}
  
  % Tables list
  {\large\bfseries\color{wistaDarkGray} Tabellen (\total{table})}\\[0.3cm]
  
  \small
  \begin{tabular}{@{}p{3cm}p{11cm}@{}}
    \toprule
    \textbf{Nr.} & \textbf{Beschreibung} \\
    \midrule
    \tablerows
    \bottomrule
  \end{tabular}
  \normalsize
  
  \vspace{1cm}
  
  % Figures list
  {\large\bfseries\color{wistaDarkGray} Abbildungen (\total{figure})}\\[0.3cm]
  
  \small
  \begin{tabular}{@{}p{3cm}p{11cm}@{}}
    \toprule
    \textbf{Nr.} & \textbf{Beschreibung} \\
    \midrule
    \figurerows
    \bottomrule
  \end{tabular}
  \normalsize
  
  \clearpage
}
\fi

% Define WISTA colors (from Kristiansen Word template)
\definecolor{wistaBlue}{RGB}{91,155,213}
\definecolor{wistaLightBlue}{RGB}{120,170,215}
\definecolor{wistaDarkGray}{RGB}{96,94,92}
\definecolor{wistaLightGray}{RGB}{128,128,128}

% Redefine bold text to use WISTA blue color
% Only redefine \textbf (inline bold), not \bfseries (used in formatting)
\let\oldtextbf\textbf
\renewcommand{\textbf}[1]{{\color{wistaBlue}\oldtextbf{#1}}}

% Section formatting (##) - number 20% bigger, lines closer, text closer to bottom line
% e.g., "2" with lines around "Einleitung"
\titleformat{\section}[display]
  {\normalfont\Large\bfseries}
  {\color{wistaBlue}\LARGE\bfseries\thesection\\\noindent\rule{\columnwidth}{1.5pt}}
  {0pt}
  {\vspace{0pt}\color{wistaDarkGray}\Large\bfseries}
  [\vspace{0pt}\noindent{\color{wistaBlue}\rule{\columnwidth}{1.5pt}}\vspace{0.05cm}]
\titlespacing*{\section}{0pt}{1.5ex plus 0.3ex minus .1ex}{0.2ex plus .1ex}

% Subsection formatting (###) - grey text with blue line under
% e.g., "2.1 Handlungsbedarf"
\renewcommand{\thesubsection}{\thesection.\arabic{subsection}}
\titleformat{\subsection}
  {\color{wistaDarkGray}\large\bfseries}
  {\color{wistaDarkGray}\large\bfseries\thesubsection}
  {0.5em}
  {\color{wistaDarkGray}\large\bfseries}
  [\vspace{0pt}\noindent{\color{wistaBlue}\rule{\columnwidth}{0.8pt}}\vspace{0.05cm}]
\titlespacing*{\subsection}{0pt}{1ex plus 0.2ex minus .1ex}{0.2ex plus .1ex}
  
% Subsubsection formatting (####) - grey text, no underline
% e.g., "2.1.1 Steigende Datenmenge"
\renewcommand{\thesubsubsection}{\thesubsection.\arabic{subsubsection}}
\titleformat{\subsubsection}
  {\color{wistaDarkGray}\normalsize\bfseries}
  {\color{wistaDarkGray}\normalsize\bfseries\thesubsubsection}
  {0.5em}
  {\color{wistaDarkGray}\normalsize\bfseries}

% Paragraph formatting (4th level) - numbered as subsection.subsubsection.paragraph
\renewcommand{\theparagraph}{\thesubsubsection.\arabic{paragraph}}
\titleformat{\paragraph}[runin]
  {\color{wistaDarkGray}\normalsize\bfseries}
  {\color{wistaBlue}\normalsize\bfseries\theparagraph}
  {0.5em}
  {\color{wistaDarkGray}\normalsize\bfseries}
  [.]
\titlespacing*{\paragraph}{0pt}{1ex plus 0.5ex minus 0.2ex}{0.5em}

% Header and footer setup with blue line under header
\pagestyle{fancy}
\fancyhf{}
\fancyhead[L]{\small\textit{\docauthors}}
\ifdraftmode
  \fancyhead[C]{\small\color{red}\textbf{ENTWURF} -- \today}
\fi
\fancyhead[R]{\small\textit{\docshorttitle}}
\fancyfoot[L]{\small\thepage}
\fancyfoot[R]{\small Statistisches Bundesamt | \docpublication\ | {\color{wistaBlue}\docpublicationissue}}
\renewcommand{\headrulewidth}{0pt}
\renewcommand{\footrulewidth}{0pt}
% Add blue line under header
\renewcommand{\headrule}{\color{wistaBlue}\hrule width\headwidth height 1pt \vskip 2pt}
\renewcommand{\headrulewidth}{1pt}

% Make all figures span both columns (full width)
\let\origfigure\figure
\let\endorigfigure\endfigure
\renewenvironment{figure}[1][htbp]{%
  \origfigure*[#1]%
}{%
  \endorigfigure*%
}

% Allow manual control of table* environment for full-width tables
% Tables by default stay in single column, but can use table* manually



% Make bibliography full-width (one column) - for biblatex
\defbibheading{bibliography}[\bibname]{%
  \onecolumn
  \section*{#1}%
  \markboth{#1}{#1}%
}



% Footnotes stay within column
\usepackage[stable]{footmisc}

% Custom footnote formatting: "|^1" in dark WISTA gray
\renewcommand{\thefootnote}{{\color{wistaDarkGray}|$^{\arabic{footnote}}$}}

% Pifont for dingbat symbols
\usepackage{pifont}

% Better table handling for two-column mode
\usepackage{array,tabularx,calc}

% Custom WISTA table format
\usepackage{caption}
\usepackage{xcolor}
\usepackage{colortbl}
\usepackage{array}

% Define custom caption format for WISTA tables
\DeclareCaptionFormat{wista}{%
  {\color{wistaLightBlue}\bfseries #1#2}\par\vspace{0.2em}%
  {\color{wistaDarkGray}\normalfont #3}\par\vspace{0.3em}%
}

% Set up table captions with custom WISTA format
\captionsetup[table]{
  position=top,
  format=wista,
  justification=justified,
  singlelinecheck=false,
  labelsep=none,
  skip=0.3em
}

% Better table spacing and white inner lines
\setlength{\tabcolsep}{6pt}
\renewcommand{\arraystretch}{1.15}

% Set white lines for tables by default, blue for first column
\AtBeginDocument{
  \arrayrulecolor{white}
  \setlength{\arrayrulewidth}{0.3pt}
}

% Define commands for colored table borders
\newcommand{\bluevrule}{\color{wistaBlue}\vrule}
\newcommand{\whitevrule}{\color{white}\vrule}

% Define WISTA table colors
\definecolor{wistaTableBg}{RGB}{240,248,255}     % Very light blue background
\definecolor{wistaHeaderBg}{RGB}{220,230,240}    % Light blue header background
\definecolor{wistaFirstCol}{RGB}{65,125,180}     % Blue for first column lines

% Custom WISTA table command
\newcommand{\wistaTableStart}{%
  \rowcolors{2}{wistaTableBg}{wistaTableBg}%
  \arrayrulecolor{wistaFirstCol}%
}

% Custom commands for table styling
\newcommand{\wistaHeaderRow}[1]{%
  \rowcolor{wistaHeaderBg}%
  \arrayrulecolor{wistaDarkGray}%
  #1 \\
  \arrayrulecolor{white}%
}

\newcommand{\wistaFirstColRule}{\arrayrulecolor{wistaFirstCol}}
\newcommand{\wistaWhiteRule}{\arrayrulecolor{white}}

% Custom bullet points with > character - 35% width, 95% height, black extra bold, half line spacing
\usepackage{enumitem}
\setlist[itemize]{label={\scalebox{0.35}[0.95]{\fontseries{bx}\selectfont >}}, leftmargin=1.5em, itemsep=5pt, topsep=3pt, parsep=0pt}

% Redefine ref to include pifont 247 (arrow) for figure/table/section references
\let\oldref\ref
\renewcommand{\ref}[1]{\oldref{#1}\raisebox{-0.2ex}{\scriptsize\color{wistaBlue}\ding{247}}}
\makeatletter
\@ifpackageloaded{caption}{}{\usepackage{caption}}
\AtBeginDocument{%
\ifdefined\contentsname
  \renewcommand*\contentsname{Inhaltsverzeichnis}
\else
  \newcommand\contentsname{Inhaltsverzeichnis}
\fi
\ifdefined\listfigurename
  \renewcommand*\listfigurename{Abbildungsverzeichnis}
\else
  \newcommand\listfigurename{Abbildungsverzeichnis}
\fi
\ifdefined\listtablename
  \renewcommand*\listtablename{Tabellenverzeichnis}
\else
  \newcommand\listtablename{Tabellenverzeichnis}
\fi
\ifdefined\figurename
  \renewcommand*\figurename{Abbildung}
\else
  \newcommand\figurename{Abbildung}
\fi
\ifdefined\tablename
  \renewcommand*\tablename{Tabelle}
\else
  \newcommand\tablename{Tabelle}
\fi
}
\@ifpackageloaded{float}{}{\usepackage{float}}
\floatstyle{ruled}
\@ifundefined{c@chapter}{\newfloat{codelisting}{h}{lop}}{\newfloat{codelisting}{h}{lop}[chapter]}
\floatname{codelisting}{Listing}
\newcommand*\listoflistings{\listof{codelisting}{Listingverzeichnis}}
\makeatother
\makeatletter
\makeatother
\makeatletter
\@ifpackageloaded{caption}{}{\usepackage{caption}}
\@ifpackageloaded{subcaption}{}{\usepackage{subcaption}}
\makeatother
\usepackage{bookmark}
\IfFileExists{xurl.sty}{\usepackage{xurl}}{} % add URL line breaks if available
\urlstyle{same}
\hypersetup{
  pdftitle={EFFIZIENTE AUSWERTUNG DES BUSINESS-TAX-PANELS MITHILFE VON R},
  pdfauthor={Oliver Hauke; Annette Kristiansen},
  pdflang={de},
  colorlinks=true,
  linkcolor={blue},
  filecolor={Maroon},
  citecolor={Blue},
  urlcolor={Blue},
  pdfcreator={LaTeX via pandoc}}


\title{EFFIZIENTE AUSWERTUNG DES BUSINESS-TAX-PANELS MITHILFE VON R}
\author{Oliver Hauke \and Annette Kristiansen}
\date{2025-12-07}
\begin{document}
\maketitle

% Custom WISTA-style title page
\begin{titlepage}
\thispagestyle{empty}
\pagestyle{empty}
\onecolumn

% Horizontal blue line at top (header line)
\noindent{\color{wistaBlue}\rule{\textwidth}{1pt}}

\vspace{-0.5cm}

% Vertical blue line on the left (6.3cm from left edge, starts below header line, stops above footer)
\begin{tikzpicture}[remember picture, overlay]
  \draw[wistaBlue, line width=2pt] 
    ([xshift=6.3cm, yshift=-4cm]current page.north west) -- 
    ([xshift=6.3cm, yshift=3.5cm]current page.south west);
\end{tikzpicture}

\vspace*{0.5cm}

% Left side - short author bios (right-aligned to the vertical line)
\begin{minipage}[t]{3.9cm}
\selectlanguage{ngerman}%
\raggedleft
\hyphenpenalty=0
\tolerance=1000
\fontsize{7}{8.5}\selectfont
\vspace{0pt}%
{\fontsize{8}{9.5}\selectfont\color{wistaBlue}\textbf{Oliver Hauke}}\\[0.15cm]
ist Mathematiker und IT-Referent im Kompetenzzentrum „Analyse und Auswertung" des Statistischen Bundesamtes. Er verantwortet die Weiterentwicklung der technischen Infrastruktur des Forschungsdatenzentrums des Statistischen Bundeamts und berät das Netzwerk empirische Steuerforschung in effizienter Datenanalyse.\\[0.5cm]

{\fontsize{8}{9.5}\selectfont\color{wistaBlue}\textbf{Annette Kristiansen}}\\[0.15cm]
ist Ökonomin und Referentin im Referat „Unternehmenssteuern" des Statistischen Bundesamtes. Sie beschäftigt sich mit der Zusammenführung der Unternehmenssteuerstatistiken und betreut Arbeiten für das Netzwerk empirische Steuerforschung. 

\end{minipage}%
\hspace{0.4cm}% Gap to vertical line (left side)
\hspace{0.4cm}% Gap from vertical line (right side)
% Right side - title, authors, keywords, abstracts
\begin{minipage}[t]{11.8cm}
\raggedright
\vspace{1.2cm}%

% Title (uppercase via command, not bold)
{\LARGE\color{wistaDarkGray}\begin{spacing}{1.3}\MakeUppercase{Effiziente Auswertung des Business-Tax-Panels mithilfe von R}\end{spacing}}

\vspace{-0.2cm}

% Blue horizontal line (underline for title)
{\color{wistaBlue}\rule{\linewidth}{0.5pt}}

\vspace{0.5cm}

% Authors inline, dark gray
{\large\color{wistaDarkGray} Oliver Hauke, Annette Kristiansen}

\vspace{0.05cm}

% Blue line (underline for authors)
{\color{wistaBlue}\rule{\linewidth}{0.5pt}}

\vspace{0.5cm}

% Keywords (pifont 247 arrow then bold label and blue text)
\hangindent=1.2em\hangafter=1
{\color{wistaBlue}\ding{247}\ {\bfseries Schlüsselwörter:} Effiziente Analysen, R-Programmierung, Forschungsdaten, Unternehmenssteuerstatistik, Business-Tax-Panel, Netzwerk empirische Steuerforschung}

\vspace{0.3cm}

% Zusammenfassung (uppercase via command)
{\color{wistaDarkGray}\bfseries\MakeUppercase{Zusammenfassung}}\\[0.15cm]
\small
Das Forschungsdatenzentrum des Statistischen Bundesamtes baut seine technische Infrastruktur gezielt aus, um der Wissenschaft erweiterte Analysemöglichkeiten amtlicher Daten bereitzustellen. Seit November 2025 steht Forschenden exklusiver Zugang zu diesen erweiterten Systemressourcen in R und Stata zur Verfügung. R-basierte Tools -- insbesondere Apache Arrow, Parquet und DuckDB -- ermöglichen es, gängige Analysen großer Forschungsdatensätze in Echtzeit auszuführen. Am Beispiel des Business-Tax-Panel (2013 bis 2019) werden die erzielbaren Effizienzgewinne bei der Auswertung umfangreicher steuerstatistischer Daten im Rahmen des Netzwerks empirische Steuerforschung demonstriert.

\vspace{0.5cm}

% English keywords and abstract (italicized)
\hangindent=1.2em\hangafter=1
{\itshape\color{wistaBlue}\ding{247}\ {\bfseries Keywords:} Efficient analysis -- R programming -- research data -- corporate tax statistics -- Business Tax Panel -- empirical tax research network}

\vspace{0.3cm}

{\itshape\color{wistaDarkGray}\bfseries\MakeUppercase{Abstract}}\\[0.15cm]
\small\itshape
The Research Data Centre at the Federal Statistical Office is strategically expanding its technical infrastructure to provide researchers with enhanced analytical capabilities for official data. Since November 2025, researchers have had exclusive access to these extended system resources in R and Stata. R-based tools -- particularly Apache Arrow, Parquet, and DuckDB -- enable real-time analysis of large research datasets. Using the Business Tax Panel (2013--2019), this paper demonstrates the achievable efficiency gains in analysing large-scale tax statistics within the empirical tax research network.

\end{minipage}

\vfill

\vspace{0.5cm}
% Footer matching document style
\noindent\makebox[\textwidth]{%
  \small 1%
  \hfill%
  \small Statistisches Bundesamt | WISTA | {\color{wistaBlue}01/2026}%
}

\end{titlepage}

% Stay in one-column mode for TOC (will be switched later)
\pagestyle{fancy}

\renewcommand*\contentsname{Inhaltsverzeichnis}
{
\hypersetup{linkcolor=}
\setcounter{tocdepth}{3}
\tableofcontents
}

\setstretch{1}
\renewcommand{\totalchars}{37 000}
\renewcommand{\totalwords}{4 489}
\renewcommand{\abstractchars}{623}
\renewcommand{\abstractwords}{85}
\renewcommand{\zusammenfassungchars}{707}
\renewcommand{\zusammenfassungwords}{81}
\renewcommand{\footnotecountnum}{16}
\renewcommand{\citationcountnum}{25}
\renewcommand{\schlusselwortercount}{6}
\renewcommand{\keywordscount}{6}
\renewcommand{\doctitle}{EFFIZIENTE AUSWERTUNG DES BUSINESS-TAX-PANELS MITHILFE VON R}
\renewcommand{\docauthors}{Oliver Hauke, Annette Kristiansen}
\renewcommand{\docpublication}{WISTA}
\renewcommand{\docpublicationissue}{01/2026}
\renewcommand{\docshorttitle}{Effiziente Analyse großer Forschungsdaten}
\renewcommand{\authoronechars}{327}
\renewcommand{\authortwochars}{256}
\newcommand{\sectionrows}{            Einleitung & \textcolor{wistaBlue}{\textbf{1 231}} & \textcolor{wistaBlue}{\textbf{11 982}} & \textcolor{wistaBlue}{\textbf{5}} & \textcolor{wistaBlue}{\textbf{8}} & --- \\
            Anwendungsfall Business-Tax-Panel & \textcolor{wistaBlue}{\textbf{125}} & \textcolor{wistaBlue}{\textbf{1 216}} & \textcolor{wistaBlue}{\textbf{1}} & \textcolor{wistaBlue}{\textbf{1}} & --- \\
            Techniche Infrastruktur & \textcolor{wistaBlue}{\textbf{145}} & \textcolor{wistaBlue}{\textbf{1 411}} & \textcolor{wistaBlue}{\textbf{1}} & \textcolor{wistaBlue}{\textbf{1}} & --- \\
            Datenverarbeitungsmethoden & \textcolor{wistaBlue}{\textbf{810}} & \textcolor{wistaBlue}{\textbf{7 884}} & \textcolor{wistaBlue}{\textbf{3}} & \textcolor{wistaBlue}{\textbf{5}} & --- \\
            Performanzvergleich & \textcolor{wistaBlue}{\textbf{999}} & \textcolor{wistaBlue}{\textbf{9 724}} & \textcolor{wistaBlue}{\textbf{4}} & \textcolor{wistaBlue}{\textbf{7}} & --- \\
            Ergebnis für das vollständige BTP & \textcolor{wistaBlue}{\textbf{84}} & \textcolor{wistaBlue}{\textbf{817}} & \textcolor{wistaBlue}{\textbf{0}} & \textcolor{wistaBlue}{\textbf{1}} & --- \\
            Fazit & \textcolor{wistaBlue}{\textbf{407}} & \textcolor{wistaBlue}{\textbf{3 961}} & \textcolor{wistaBlue}{\textbf{2}} & \textcolor{wistaBlue}{\textbf{3}} & --- \\
}
\newcommand{\tablerows}{            \hyperref[tbl-stba-rserver]{Tabelle~1} & StBA R-Server Ausbau (Anfang 2025) \\
            \hyperref[tbl-fdz-server]{Tabelle~2} & FDZ OnSite-Server Spezifikationen \\
            \hyperref[tbl-pkg-versions]{Tabelle~3} & Versionen verfügbarer R-Packages \\
            \hyperref[tbl-baseline]{Tabelle~4} & Baselines in SAS und Stata (anhand 1\% des simuliertes BTP) \\
            \hyperref[tbl-benchmark-small]{Tabelle~5} & Performanz Messungen für die Fallzahlberechnung anhand 1\% des BTP \\
            \hyperref[tbl-benchmark-regr-small]{Tabelle~6} & Performanz Messungen für die Regressionsberechnung anhand 1\% des BTP \\
            \hyperref[tbl-benchmark]{Tabelle~7} & Performanz Messungen für die Fallzahlberechnung anhand 10\% des BTP \\
            \hyperref[tbl-benchmark-regr]{Tabelle~8} & Performanz Messungen für die Regressionsberechnung anhand 10\% des BTP \\
            \hyperref[tbl-bm-cnt-small]{Tabelle~9} & Benchmarkergebnisse für Fallzahlen (kleine Stichproben des BTP) \\
            \hyperref[tbl-bm-reg-small]{Tabelle~10} & Benchmarkergebnisse für Regression (kleine Stichproben des BTP) \\
}
\newcommand{\figurerows}{            \hyperref[fig-optimization-strategies]{Abbildung~1} & Strategien zur Performanz-Optimierung im FDZ-Bund \\
}
\newcommand{\schlusselworterlist}{\textcolor{wistaBlue}{\ding{247}} Effiziente Analysen\\
\textcolor{wistaBlue}{\ding{247}} R-Programmierung\\
\textcolor{wistaBlue}{\ding{247}} Forschungsdaten\\
\textcolor{wistaBlue}{\ding{247}} Unternehmenssteuerstatistik\\
\textcolor{wistaBlue}{\ding{247}} Business-Tax-Panel\\
\textcolor{wistaBlue}{\ding{247}} Netzwerk empirische Steuerforschung\\
}
\newcommand{\keywordslist}{\textcolor{wistaBlue}{\ding{247}} Efficient analysis\\
\textcolor{wistaBlue}{\ding{247}} R programming\\
\textcolor{wistaBlue}{\ding{247}} research data\\
\textcolor{wistaBlue}{\ding{247}} corporate tax statistics\\
\textcolor{wistaBlue}{\ding{247}} Business Tax Panel\\
\textcolor{wistaBlue}{\ding{247}} empirical tax research network\\
}
\ifdraftmode
\metadatapage
\fi

% Table of Contents will be inserted here by Quarto
\clearpage
% Switch to two-column mode for main document
\twocolumn

\section{Einleitung}\label{sec-einleitung}

\begin{enumerate}
\def\labelenumi{\arabic{enumi}.}
\tightlist
\item
  Steigende Datenmengen, höhere Bedarfe der Öffentlichkeit an Aktualität
  und Detailliertheit -\textgreater{} Modernisierung der
  Datenverarbeitungsinfastruktur nötig um mithalten zu können
\item
  Zugespitz in FDZ. Forschende fordern mehr umfangreiche Daten (BMF
  Gutachten) und die technischen Möglichkeiten sie auszuwerten (SVR
  Gutachten)
\item
  Gründung NeSt, neue umfangreiche Datenquelle BTP mit hohem
  Analysepotenzial (S. \textbf{?@sec-btp})
\item
  Große Schritte in 2025, erweiterung der R-Server. NeSt auch
  modernisierung der FDZ Infrastruktur geführt. Technisches Fundament
  (S. Kapitel~\ref{sec-infrastruktur})
\item
  Vorteile von R Open Source, vielzahl R-Packages, aktive Community,
  wachsende Verbreitung in OSS, natürliche Wahl R. bieten Effizienz,
  hohe Nutzerfreundlichkeit (s Kapitel~\ref{sec-methoden})
\item
  Frage nach besten möglichkeiten, die Anforderungen XYZ bedient (s.
  Kapitel~\ref{sec-benchmark})
\item
  zeigen, dass die Methoden sogar das ganze BTP sehr effizient auswerten
  können, in Sek, obwohl es den verfügbaren Arbeitsspeicher übersteigt
  (s. Kapitel~\ref{sec-ergebnis}). Ableitung praktischer Implikationen.
\end{enumerate}

Insbesondere die Wissenschaft erwartet, das die
\emph{Forschungsdatenzentrum des statistischen Bundesamts (hier kurz
FDZ) und dem Forschungsdatenzentrum der Statistischen Ämter der Länder}
zeitnah auch umfangreiche Datensätze für ausgiebigere wissenschaftliche
Auswertungen zur Verfügung stellt. In dem FDZ haben Forschende die
Möglichkeit, eigene Analysen mit amtlicher Mikrodaten durchführen. Der
Zugangsweg ``Onsite-Nutzung'' ermöglicht den Zugriff zu formal
anonymisierten Einzeldaten ermöglichen. Forschende können selbst ein
\emph{Gastwissenschaftsarbeitsplatz} (GWAP) in den Liegenschaften des
FDZ nutzen oder via \emph{kontrollierter Datenfernverarbeitung} (KDFV)
können entwickelte Auswertungsprogramme durch das FDZ ausgeführt werden.
Jedoch hatte die Wissenschaft in \textcite{bmf2020dateninfrastruktur}
Kritik am Datenbestand für empirische Steuerforschung und dessen
Aktualität geäußert.In \textcite{svr2023jahresgutachten} wurde darüber
hinaus die technisch veraltete Infrastruktur des FDZ bemängelt.

Um diese Situation zu verbessern, hat das Bundesfinanzministerium im
Jahre 2022 das \emph{``Netzwerk empirischer Steuerforschung''} (NeSt)
gegründet. \footnote{https://www.bundesfinanzministerium.de/Web/DE/Themen/Steuern/Steuerliche\_Themengebiete/NeSt/netzwerk-fuer-empirische-steuerforschung.html}
Vorangetrieben durch das NeSt konnte das Statistische Bundesamt bereits
erfolgreich den Datenbestand erweitern. Insbesondere steht seit 2023 das
\emph{Business-Tax-Panel} (BTP) zur Verfügung und wird bereits erneuert
(in \textbf{?@sec-btp} beschrieben). Es veknüpft verschiedene
Unternehmensstatistiken und birgt mit fast 3 000 Merkmalen einerseits
ein großes Analysepotenzial, andererseits hohe technische Anforderungen
an die Analysewerkzeuge.\footnote{https://www.forschungsdatenzentrum.de/de/steuern/btp}

\subsection{Voranschreitende
Innovation}\label{voranschreitende-innovation}

Anforderungen: Performanz auf neuem Stand der Technik, schnelle
Ergebnisse, idealenfalls in Echtzeit (d.h. wenigen Sekunden). Leichte
Nutzung / geringe Entwickleraufwände (im FDZ für viele Fachrichtungen
zugänglich), moderne Interfaces, möglichst leiche Nutzung, wenig
Schulungsbedarf. leicht einrichtbar, wenig Abhängigkeiten zu vielen
Tools, die sich ständig weiterentwicklen, sodass hohe Wartungsaufwände
entstehen. stabil, optimale Nutzung vorhandener GWAP Zeit und KDFV
Zeitraum.ressourcensparend (RAM Problematik)

Um diesen Herausforderungen zu bewältigen, hat das \emph{Statstische
Bundesamt} in 2025 wesentliche Meilensteine in der Modernisierung der
technischen Analyseinfrastruktur erreicht (dargestellt in
Kapitel~\ref{sec-infrastruktur}). Anfang 2025 wurden die R-Server für
Auswertungen und Datenverabeitungen in der Statistikproduktion
modernisiert und stark erweitert. Seit November 2025 stehen auch dem FDZ
neue R- und Stata-Analyseserver (sog. FDZ OnSite-Server) für
wissenschaftliche Auswertungen zu Verfügung.

Um die technischen Mittel bereitzustellen, hat eine weitere Initiative
des NeSt die Moderniserung IT-Infrastruktur des FDZ vorangetrieben.
\textbf{Dadurch bietet das FDZ seit November 2025 neue Server mit großen
Systemkapazitäten für R und Stata für wissenschaftliche Auswertungen
an}. Erstmalig können die vollen Kapazitäten neben der KDFV auch an den
GWAPs genutzt werden, um große Datensätze im Vollmaterial auszuwerten.
Wir werden in den FDZ OnSite-Servern demonstrieren, dass viele
Auswertungsbedarfe das BTP im Vollmaterial innerhalb von Sekunden
beantwortet werden können. Da die Datenmenge des BTP den verfügbaren
Arbeitsspeicher weiterhin übertrifft, sind geeigneten
Programmiermethoden notwendig (untersucht in
Kapitel~\ref{sec-ergebnis}).

Neben den bisherigen Hauptwerkzeuge für Datenverarbeitung sind für die
Statistikproduktion aktuell SAS und Stata in dem FDZ, findet R als
Lingua Franca in der Statistik, Methodologie und Datenwissenschaft und
somit auf natürlich weise Einzug in die amtliche Statistik.\footnote{Wachstum
  der uRos s. https://r-project.ro/conference2025.html\#Presentations}

Als Open Source Software mit einer weltweit aktiven Community, bietet R
bietet Zugang zu sehr vielen Funktionen der Datenverarbeitung (in Form
von sog. \emph{R-Paketen}), die ein breites Spektrum an
Nutzerfreundlichkeit und Performanz bieten. Die gleiche Datenvearbeitung
kann abhängig von der Software-Nutzung, einige Stunden dauern oder in
wenigen Sekunden fertigs ein. Nur die richtige Nutzung geeigneter
R-Pakete ermöglicht es, Datenverarbeitungsbedarfe anhand großer
Datenmengen in Sekunden nachvollziehbar auszuführen (beschrieben in
Kapitel~\ref{sec-methoden}). Dadurch stellt sich die Frage welche
Programmiermethoden geeignet sind, die neuen Datenverarbeitungsbedarfe
mit der spezifischen neuen Infrastruktur zu bedienen (diskutiert in
Kapitel~\ref{sec-benchmark}). Durch die Beantwortung können für den
Statistikerstellungsprozesses Werkzeuge abgeleitet werden, die
Informationsbedarfe der Öffentlichkeit schnell anhand wachsender
Datenmengen zu bedienen. Die Bewertung muss sorgfältig durchgeführt
werden, da sich die betrachteten R-Pakete schnell weiterentwicklung,
meist mit mehreren großen Releases innerhalb eines Jahres, aber die
organisatorischen Strukturen nicht erlauben, jede neue Version
einzurichten.

\begin{verbatim}
-   wechselnde Nutzende (Forscher) unterschiedlicher Fachbereiche mit unterschiedlichen Analytischen und technischen Kenntnissen
\end{verbatim}

\section{Anwendungsfall
Business-Tax-Panel}\label{anwendungsfall-business-tax-panel}

\subsection{Datengrundlage}\label{datengrundlage}

\emph{Das Business-Tax-Panel kombiniert die Angaben aus verschiedenen
Ertrags- sowie Umsatzsteuerstatistiken im Quer- und Längsschnitt.
Verfügbar sind Informationen aus den folgenden Statistiken: Gewerbe-,
Körperschaftsteuer, Umsatzsteuer-Voranmeldung und -Veranlagung,
Statistik über die Personengesellschaften und Gemeinschaften,
Einnahmenüberschussrechnung und einige Merkmale aus dem Statistischen
Unternehmensregister. Insgesamt deckt der Forschungsdatensatz eine
Vielzahl an Analysemöglichkeiten, wie zum Beispiel Analyse von
Anpassungsreaktionen oder Verteilungsanalysen, die für die
evidenzbasierte steuerpolitische Entscheidungen nötig sind, ab.}

\emph{Für die vielfältigen Analysewünsche der Forschenden, der damit
einhergenden Dokumentation des Datensatzes in den Metadatenreports
(!!hier Fußnote für MDR - siehe Deskriptive Ecktaten 1.) des FDZ, und
den Auswertungsbedarf des Fachbereiches ist der Datensatz möglichst
wenig beschränkt worden und dadurch ressourcenintensive. Anstelle einer
Stichprobe wurden die zugrundeliegenden Vollerhebungen verwendet. Die
erste Version des BTP (2013 bis 2019) besteht aus ca. 67 Millionen
Beobachtungen. Mit der Aktualisierung um das Berichtsjahr 2020 wird die
Zahl der Beobachtungen sich auf ca. 78 Millionen erhöhen. Auch der
Merkmalsumfang von zurzeit über 2.700 Merkmalen wird durch Erweiterung
von Berichtsjahren steigen.}

\emph{Bereits aktuell und insbesondere durch die regelmäßige Erweiterung
um Bereichtsjahre steht der Fachbereich sowie das FDZ vor der
Herausforderung den hohe Bedarf an Speicherkapazität, Arbeitsspeicher
und Prozessorleistung zu decken. Daher wird das BTP als Anwendungsfall
herangezogen}

\subsection{Auswertungsbedarfe}\label{auswertungsbedarfe}

\emph{Um den Auswertungsbedarf des Fachbereichs und des FDZ abzubilden
wird die Tabelle zur Häufigkeitsverteilung der Beobachtungen je Jahr und
Steuerstatistik (FDZ, 2024 --\textgreater{} IN QULLEN aufnehmen!)
verwendet. Im Fachbereich ergibt sich zusätzlich der Bedarf die Daten je
nach Statistik und Tabelle zu bearbeiten. Auch die Forschenden
bearbeiten die Daten um sie für die jeweilige Analyse aufzubereiten.
Daher wird die Datenbearbeitung und anschließende Analyse als zweiter
Anwendungsbedarf getestet.}

\emph{Die Körperschaftsteuerstatistik zeigt über die Jahre eine Zunahme
der Steuerpflichtigen mit Verlustvortrag im jeweiligen Berichtsjahr
(GENESIS Code 73211-0002). Für den zweiten Anwendungsfall wird das
Vorliegen eines Verlustvortrages im jeweiligen Berichtsjahr als versucht
anhand verschiedener Einflussfaktoren zu erklären. Anzunehmen ist, dass
Personalkosten den Verlustvortrag erhöhen können, daher wird die Zahl
der Sozialversicherungspflichtig Beschäftigen sowie die Summe der Löhne
und Gehälter als erklärende Variable verwendet. Der jeweilige
Gesamtbetrag der Einkünfte und die abgezogenen Spenden können ebenfalls
den Verlustvortrag beeinflussen. Um den Internationalisierungsgrad des
Unternehmens zu erfassen wird ein Dummy für aufgenommen der Anzeigt, ob
das Unternehmen auslandskontrolliert ist.}

\subsubsection{Deskriptive Eckdaten}\label{deskriptive-eckdaten}

\begin{enumerate}
\def\labelenumi{\arabic{enumi}.}
\tightlist
\item
  deskriptive Eckdaten, für BTP als Panel im Quer- und Längsschnitt:
  Fallzahl pro Statistik und Jahr (s.MDR Produkt S. 9f
  https://www.forschungsdatenzentrum.de/sites/default/files/btp\_2013-2019\_on-site\_mdr-prod.pdf)
\end{enumerate}

Kennzahlen um inhaltlose Testdaten zu generieren werden folgende
Eckdaten benötigt und im nächsten BTP berücksichtigt:

\begin{itemize}
\tightlist
\item
  Befüllung
\item
  Statistische Eckdaten univariate (Quantile),
\item
  ggf. multivariate (gemeinsame Verteilung)
\end{itemize}

Exemplarisch, neues BTP beinhaltet dies. Regelmäßige Updates
beschleunigen.

\subsubsection{Regressionen}\label{regressionen}

\begin{enumerate}
\def\labelenumi{\arabic{enumi}.}
\setcounter{enumi}{1}
\tightlist
\item
  Verlustvorträge\footnote{Vortrag Annette 2024 JK} (s. Vortrag NeSt)
  als Indikator für \ldots{} und Einflussgrößen auf Verlustvorträge, für
  Vermutung von Einfluss der Unternehmensgröße (großes Gesamtvolumen,
  Löhne, Spenden, tätige Personen; Auslandskontrolle oder WZ) auf
  positive Verlustvorträge.
\end{enumerate}

\subsection{Technische
Spezifikationen}\label{technische-spezifikationen}

\emph{AK: Habe ich im Bereich Datengrundlage integriert. }

\emph{Für die vielfältigen Analysewünsche der Forschenden, den
Auswertungsbedarf des Fachbereiches bspw. für Sonderauswertungen und der
Erstellung des Metadatenreports des Forschungsdatenzentrums ist der
Forschungsdatensatz konzipiert. Für die Abdeckung aller Anforderungen
ist der Datensatz möglichst wenig beschränkt worden und dadurch
ressourcenintensive. Der damit einhergehende hohe Bedarf an
Speicherkapazität, Arbeitsspeicher und Prozessorleistung stellt sowohl
im Fachbereich als auch im Forschungsdatenzentrum eine Herausforderung
dar.}

\emph{Für die vielfältigen Analysewünsche der Forschenden, den
Auswertungsbedarf des Fachbereiches bspw. für Sonderauswertungen und der
Erstellung des Metadatenreports des Forschungsdatenzentrums ist der
Forschungsdatensatz konzipiert. Für die Abdeckung aller Anforderungen
ist der Datensatz möglichst wenig beschränkt worden und dadurch
ressourcenintensive. Der damit einhergehende hohe Bedarf an
Speicherkapazität, Arbeitsspeicher und Prozessorleistung stellt sowohl
im Fachbereich als auch im Forschungsdatenzentrum eine Herausforderung
dar.}

\begin{itemize}
\item
  Technische Herausforderung:

  \begin{itemize}
  \tightlist
  \item
    effiziente und verständliche Speicherung
  \item
    Herausfiltern großer Datenmengen irrelevanter Merkmale
  \end{itemize}
\item
  Umfang über verfügbaren Arbeitsspeicher
\item
  Besonders breit mit 3000 Merkmalen, 69 Mio. Beobachtung, sodass
  zeilenbasierte Formate problematisch sind.
\item
  Technische Herausforderung:

  \begin{itemize}
  \tightlist
  \item
    effiziente und verständliche Speicherung
  \item
    Herausfiltern großer Datenmengen irrelevanter Merkmale
  \end{itemize}
\end{itemize}

\section{Ausbau der Analyseinfrastruktur}\label{sec-infrastruktur}

Im Jahr 2025 konnten mit der Bereitstellung neuer Analyseserver
wesentliche Fortschritte beim Ausbau der analytischen Infrastruktur des
Statistischen Bundesamts und des FDZ erzielt werden.

Die Einführung der neuen Systeme stellt einen bedeutenden Schritt dar,
da der Ausbau der Infrastruktur in den formalen Strukturen der
Bundesverwaltung ein mehrjähriges Unterfangen ist. Auch Änderungen von
Softwareversionen in der abgeschotteten Infrastruktur können
organisatorisch aufwändig sein, sodass für die folgenden Darstellungen
nicht immer die neuesten Versionen der betrachteten Software zur
Verfügung stehen --- ein Umstand, der die realistischen
Rahmenbedingungen der laufenden Statistikproduktion widerspiegelt. Um
dennoch langfristig leistungsstark zu bleiben, wurde die neue
Infrastruktur von Beginn an zukunftsorientiert konzipiert und bietet
moderne Software auf einer deutlich erweiterten Hardwarebasis.

\subsection{Ausgangslage}\label{ausgangslage}

Das Hauptwerkzeug für die Statistikproduktion ist die kommerzielle
Analysesoftware \textbf{SAS 9.4}, für das dedizierte Server
bereitgestellt werden.\footnote{Da SAS bereits langjährig eingesetzt
  wird, ist für die kommenden Performanzvergleiche zu beachten, dass die
  SAS-Server nicht über die neuste Hardware verfügen. SAS wurde in 2025
  auf Version 9.4M8 aktualisiert. Die neuste Version ist SAS 9.4M9. S.
  https://documentation.sas.com/doc/en/pgmsascdc/9.4\_3.5/whatsdiff/upgradeswhatsnew94.html}
SAS 9.4 verarbeitet Daten überwiegend zeilenorientiert und belastet den
Arbeitsspeicher daher nur gering --- ein entscheidender Vorteil in der
amtlichen Statistik, da auch große Datenmengen auf mit geringeren
Hardwareanforderungen verarbeitet werden können. Für viele
Anwendungsfälle ist jedoch die spaltenbasierte Datenverarbeitung
effizienter. In R sind neuere Verfahren verfügbar, die moderne
Prozessor- und Arbeitsspeicher gezielt ausnutzen, können die Laufzeiten
gegenüber SAS erheblich verkürzen (siehe Kapitel~\ref{sec-methoden} und
Kapitel~\ref{sec-benchmark}).

Im FDZ ist -- neben SAS -- die kommerzielle Software \textbf{Stata} das
am stärksten nachgefragte Auswertungswerkzeug. \textbf{Stata} wird in
vielen wissenschaftlichen Fachrichtungen genutzt, da sie ohne
tiefgreifende Programmierkenntnisse eingesetzt werden kann, gleichzeitig
aber spezialisierte Verfahren, insbesondere für ökonometrische Analysen,
bereitstellt. Bei großen Datenmengen stößt Stata jedoch an seine
Grenzen, da Daten vollständig in den Arbeitsspeicher geladen werden
müssen, und nur ausgewählte Operationen erlauben die parallelisierte
Datenverarbeitung. Durch die manuelle \emph{Variablenauswahl} lassen
sich größere Datensätze bearbeiten, der Speicherbedarf bleibt aber
begrenzend. Bisher waren die GWAP-Arbeitsplätze auf die
Single-Core-Version mit begrenztem RAM beschränkt und damit für viele
Produkte des FDZ unzureichend, während leistungsfähigere Server vor der
Bereitstellung des OnSite-Server nur in der KDFV verfügbar waren.

\subsection{R-Server des Statistischen
Bundesamts}\label{r-server-des-statistischen-bundesamts}

Seit Jahren steht im Statistischen Bundesamt ein Server für die
Open-Source \textbf{Statistiksoftware R} und \textbf{Python} zur
Verfügung, der für die Datenverarbeitung in der Statistikproduktion
sowie im KDFV des FDZ genutzt wird. R wird aufgrund seiner großen
internationalen Community, der Vielzahl spezialisierter Pakete für
Datenanalyse, spaltenbasierte Verarbeitung und Visualisierung sowie der
schnellen Weiterentwicklung zunehmend in der amtlichen Statistik
eingesetzt.

Seit Anfang 2025 konnte der R- und Python-Server signifikant erweitert
werden. Verwendet wird die kommerziell unterstützte \textbf{Posit
Workbench}, um die Funktionen Projekt-Sharing, Git-Integration,
Positron-Features bereitstellen zu können. Zusätzlich wurden die
Systemressourcen deutlich ausgebaut, wie die
Tabelle~\ref{tbl-stba-rserver} zeigt.

\begin{table*}

\caption{\label{tbl-stba-rserver}StBA R-Server Ausbau (Anfang 2025)}

\centering{

[!htb]

\centering\begingroup\fontsize{10}{12}\selectfont
\begin{tabular}{|>{}l|||>{}l|||>{}l|}
\hline
\cellcolor[HTML]{dce6f0}{\textcolor[HTML]{404040}{\textbf{Komponente}}} & \cellcolor[HTML]{dce6f0}{\textcolor[HTML]{404040}{\textbf{Vorher}}} & \cellcolor[HTML]{dce6f0}{\textcolor[HTML]{404040}{\textbf{Nachher}}}\\
\hline
\cellcolor{white}{\textcolor{black}{Anzahl Server}} & \cellcolor[HTML]{f0f8ff}{1} & \cellcolor[HTML]{f0f8ff}{10}\\
\cellcolor{white}{\textcolor{black}{Arbeitsspeicher (jeweils)}} & \cellcolor[HTML]{f0f8ff}{700 GB} & \cellcolor[HTML]{f0f8ff}{1 TB}\\
\cellcolor{white}{\textcolor{black}{Kernanzahl (jeweils)}} & \cellcolor[HTML]{f0f8ff}{72} & \cellcolor[HTML]{f0f8ff}{96}\\
\cellcolor{white}{\textcolor{black}{Core Speed (jeweils)}} & \cellcolor[HTML]{f0f8ff}{TBA} & \cellcolor[HTML]{f0f8ff}{TBA}\\
\cellcolor{white}{\textcolor{black}{GPUs (jeweils)}} & \cellcolor[HTML]{f0f8ff}{keine} & \cellcolor[HTML]{f0f8ff}{\makecell[l]{2× A6000\\(je 48 GB)}}\\
\hline
\end{tabular}
\endgroup{}

}

\end{table*}%

Die neue Oberfläche und IDE erhöhen die Nutzerfreundlichkeit, die
neuesten R-Versionen stehen bereit, und Pakete können schnell
installiert werden. So können moderne Analyse- und
Visualisierungsmethoden effizient genutzt werden -- sodass aktuell über
200 aktive Nutzende das System nutzen.

\subsection{OnSite-Server des FDZ}\label{onsite-server-des-fdz}

Seit November 2025 steht der sogenannte OnSite-Server exklusiv im FDZ
für alle OnSite-Nutzungen bereit. Er bietet jeweils dedizierte Stata-
und R-Server an den GWAP-Arbeitsplätzen sowie via KDFV mit gleichermaßen
umfangreichen Analysekapazitäten, wie in Tabelle~\ref{tbl-fdz-server}
aufgeführt. Alle Analyseserver sind über schnelle Direktverbindungen an
den Netzwerkspeicher angebunden, sodass R und Stata nahtlos kombiniert
werden können.

\begin{table*}

\caption{\label{tbl-fdz-server}FDZ OnSite-Server Spezifikationen}

\centering{

[!htb]

\centering\begingroup\fontsize{10}{12}\selectfont
\begin{tabular}{|>{}l|||>{}l|}
\hline
\cellcolor[HTML]{dce6f0}{\textcolor[HTML]{404040}{\textbf{Komponente}}} & \cellcolor[HTML]{dce6f0}{\textcolor[HTML]{404040}{\textbf{Umfang}}}\\
\hline
\cellcolor{white}{\textcolor{black}{Stata-Server}} & \cellcolor[HTML]{f0f8ff}{3}\\
\cellcolor{white}{\textcolor{black}{R-Server}} & \cellcolor[HTML]{f0f8ff}{6}\\
\cellcolor{white}{\textcolor{black}{Arbeitsspeicher (jeweils)}} & \cellcolor[HTML]{f0f8ff}{512 GB}\\
\cellcolor{white}{\textcolor{black}{Kernanzahl (jeweils)}} & \cellcolor[HTML]{f0f8ff}{16}\\
\cellcolor{white}{\textcolor{black}{CPU-Geschwindigkeit (jeweils)}} & \cellcolor[HTML]{f0f8ff}{3.1 GHz}\\
\hline
\end{tabular}
\endgroup{}

}

\end{table*}%

Stata/MP ist nun einheitlich auf leistungsstarken Multi-Core-Servern
verfügbar. Zudem steht erstmals R für OnSite-Nutzungen bereit: Aufbau
und Softwarebasis orientieren sich am zentralen R-Server des
Statistischen Bundesamts, sind jedoch dediziert für FDZ-Anwendungen.

Die R-Server unterstützen moderne, speichereffiziente Verfahren, mit
denen auch Datenmengen oberhalb des Arbeitsspeichers flexibel, stabil
und schneller verarbeitet werden können -- Details dazu siehe
Kapitel~\ref{sec-methoden}.

\section{Datenverarbeitungsmethoden}\label{sec-methoden}

Wir betrachten in diesem Artikel die Performanz von sog.
\emph{Datenverarbeitungsmethoden}, d.h. Softwarelösungen, die Daten
einlesen, modifzierern, aggregieren, auswerten oder
analysieren.\footnote{Das Schreiben von Daten wird in unserer
  Untersuchung vernachlässigt, da es sich bei den Veröffentlichungen des
  Statistischen Bundesamtes oder den Ergebnissen der Wissenschaft im FDZ
  um kleine Tabellen handelt, die meist keine große Performanz
  erfordern. Dennoch bieten die in in Kapitel~\ref{sec-highperf}
  vorstellten Lösungen auch hohe Performanz im Schreiben von Daten.
  Aufgrund der externen Abhängigkeiten verzichten wir auf die
  betrachtung von Datenbanken und Spark.}

\subsection{Anforderungen an
Datenverarbeitungsmethoden}\label{anforderungen-an-datenverarbeitungsmethoden}

Die wichtigsten Anforderungen bedient R anhand der folgenden Funktionen.
- hohe Usablity - hohe Performanz

Optimierung der Gesamtaufwände (incl.~Entwicklerzeit), nicht nur
Laufzeit (\textcite{gomolka2021largedata}).

\subsection{Klassische Methoden}\label{klassische-methoden}

R-Base ist zweckdienlich für Einsteiger, abhängigkeitsfreie
Infrastrukturen, jedoch weder in Usablity, noch Performanz
empfehlenswert. Basieren auf CSV-Format. R Objekte können auch effizient
binär abgespeichert werden rds, qs.

Hervorragende Ergänzungen zu R-Base:

\begin{itemize}
\tightlist
\item
  \textbf{data.table}. greift die R-Basissyntax auf, aber ist das
  optimal schnelle Package Basisfunktionen in R in Userexperience und
  Performanz nicht empfehlenswert

  \begin{itemize}
  \tightlist
  \item
    \{data.table\} \autocite{R-data.table} in den meisten Fällen die
    performanteste klassische DVM. An die Syntax an base-r angesiedelt
    und nur hier verwendet, spezielles Wissen notwendig.
  \item
    \textbf{Vorteile:} Sehr schnell für In-Memory-Operationen,
    ausgereift, kompakte Syntax In teilen der Community sehr beliebt
  \item
    \textbf{Nachteile:} Alle Daten müssen in RAM passen, steile
    Lernkurve, keine Out-of-Core-Unterstützung.
  \item
    \textbf{Empfehlung:} Gute Alternative wenn Datensatz in RAM passt
    und keine Parquet-Integration benötigt wird
  \end{itemize}
\item
  \textbf{tidyverse} (readr, dplyr und weitere): \ldots{}

  \begin{itemize}
  \tightlist
  \item
    Tidyverse \autocite{R-tidyverse} ist ein Ökosystem von Packages zur
    Datenverarbeitung, das für besonders hohe Lesbarkeit des Codes
    bekannt ist. Dazu dient eine eigenen Programmiersyntax, sog.
    \emph{\{dplyr\}-Syntax} \autocite{R-dplyr}, die den Anspruch hat
    menschlicher Sprache abzubilden. Das R-Package \{dplyr\} ist die
    Datenverarbeitungsgrammatik innerhalb Tidyverse, die anhand eines
    Beispiels in der folgenden Tabelle veranschaulicht wird
  \item
    Diese Syntax ist bei einem Teil der R-Community so beliebt, dass sie
    vielen anderen Packages ergänzt wurde. Bspw. für \{data.table\}
    setzen dies die Packages \{dtplyr\} \autocite{R-dtplyr} oder
    \{tidytable\} \autocite{R-tidytable} um Wie im Falle von \{dtplyr\}
    kann die Übersetzung zu \{dplyr\} jedoch die Performanz
    verschlechtern.\footnote{S.
      https://www.r-bloggers.com/2024/04/the-truth-about-tidy-wrappers-3/}
  \item
    nutzerfreundlich, einfach zu lernen / nutzen
  \item
    vroom
  \end{itemize}
\item
  \textbf{collapse}
\end{itemize}

{[}TABELLE DPLYR BEISPIEL{]}

\subsection{Hoch-performante Methoden in R}\label{sec-highperf}

Direkte Performanzgewinne

\begin{itemize}
\tightlist
\item
  Variablenselektion, e.g in BBK: Bot Beschleuniging von 39.9\% in R und
  33.9\% in Stata
\end{itemize}

Ab mittleren Datenmengen (relativ zu CPU und RAM) sind ausgewählte
\textbf{Best Practices} in R:

\begin{itemize}
\tightlist
\item
  \emph{Variablenselektion}: Beim Einlesen von DatenAuswahl der
  benötigten Variablen in der Einlesefunktion ist wesentliche
  Best-Practices (besonders für klassiche Methoden, da dies das Risiko
  für RAM-Abbrüche und Wartezeiten deutlich reduziert). Das Einlesen
  zunächst aller Daten (ggf. mit Variablenselektion zu einem späteren
  Zeitpunkt) verursacht unnötige Wartezeiten, Ressourcenbelegung und
  kann zu Abbrüchen führen.
\item
  \emph{Datenaufteilung}: Wenn der Arbeitsspeicher nicht ausreicht, um
  die benötigten Daten einzulesen, bietet eine Aufteilung auf mehrere
  RAM-Gerechte Datenpakete helfen. Für klassische Methoden kombiniert
  mit der mehrfachen Variablenselektion. Die folgenden hoch-performanten
  Techniken, kommen auch direkt mehreren Dateien nutzen, ohne manuell
  die Datensätze zusammenzuführen.
\end{itemize}

z.\,B. durch Lazy-Computation\footnote{Ohne das Nutzende ihren Code
  verändern müssen, speichert die Lazy-Computation geplante
  Datenverarbeitungsschritte im Hintergrund, ohne sie sofort auszufüren.
  Erst wenn ein konkretes Ergebnis abgefragt wird, führt das System die
  Berechnung durch. Dadurch werden nur die tatsächlich benötigten Daten
  verarbeitet und unnötige Schritte vermieden -- ohne dass der Nutzende
  seinen Code ändern muss.} -- ohne dass eine manuelle Variablenauswahl
erforderlich ist (siehe Kapitel~\ref{sec-methoden}).

\begin{itemize}
\tightlist
\item
  Parquet

  \begin{itemize}
  \tightlist
  \item
    R zeigt große Performanzvorteile gegenüber Stata für FDZ Aufgaben:
    ``Note that using parquet and R reduces the computing time compared
    to the optimised Stata approach by 65.1\% from 77.6 to 27.1
    seconds'' (\textcite{gomolka2021largedata})
  \item
    Die u.g. Parquet-Dateien werden bereits standardmäßig im FDSZ
    angeboten.
  \end{itemize}
\item
  arrow

  \begin{itemize}
  \tightlist
  \item
    c++ basiert
  \end{itemize}
\item
  duckdb

  \begin{itemize}
  \tightlist
  \item
    \ldots{}
  \end{itemize}
\item
  polars

  \begin{itemize}
  \tightlist
  \item
    rust-basiert
  \end{itemize}
\end{itemize}

Diverse benchmarks kommen zu unterschiedlichen Ergebnissen, daher ist
die untersuchung in eigener Infrastruktur mir relevanten Use Cases
wichtig und wird im folgenden Kapitel~\ref{sec-benchmark}
durchgeführt.\footnote{Benchmark von Data.table:
  https://cran.r-project.org/web/packages/data.table/vignettes/datatable-benchmarking.html}\footnote{\ldots{}}

Wir sprechen in diesem Artikel vereinfacht von einer
\emph{hoch-performanten Methode}, sobald die betrachtete
Datenverarbeitungsmethode in der Lage ist, Datenmengen überhalb des
verfügbaren Arbeitsspeichers flexibel und effizient zu nutzen. SAS
erfüllt den Großteil dieser Definition -- lediglich die Effizienz lässt
sich deutlich steigern.

Um Datenauswertung auf höhere Datenmengen skalieren zu können ist ein
Paradigmenwechsel von der ``\emph{vollständigen Verarbeitung im RAM''}
zu der \emph{Auslagerung aus dem RAM} notwendig. Vollständig im RAM sind
Daten zwar sehr schnell und einfach zu nutzen, jedoch beansprucht der
Transport in den RAM (Einlesen) einen Hauptteil der Laufzeit und der
verfügbare RAM reicht i.d.R. nicht für größere Datenmengen. Die
nachfolgenden Technologien sind auch in der Lage Daten überhalb des RAMs
zu verarbeiten. Durch moderne Hardware und Software kann dies sogar sehr
schnell erfolgen.

\subsubsection{Apache Parquet}\label{apache-parquet}

\begin{itemize}
\tightlist
\item
  Datenformat, spaltenbasiert, komprimiert, \ldots{}
\item
  tba - s. Vortrag
\end{itemize}

\subsubsection{Apache Arrow}\label{apache-arrow}

\begin{itemize}
\tightlist
\item
  R-package \autocite{R-arrow}, C++ basiert, erlaubt Parallelisierung,
  Datenaufteilung, Lazy-Computations, Datenformat im RAM passend zu
  Parquet
\item
  tba - s. Vortrag
\end{itemize}

\subsubsection{DuckDB}\label{duckdb}

\begin{itemize}
\tightlist
\item
  R-Package \autocite{R-duckdb}, in-memory DB, \ldots{}
\item
  tba - s. Vortrag
\end{itemize}

\subsubsection{Polars}\label{polars}

\begin{itemize}
\tightlist
\item
  an Python-angelehnte Syntax \autocite{R-polars}
\item
  tba
\item
  Benchmark des Polars Entwicklerteams \footnote{https://ddotta.github.io/cookbook-rpolars/benchmarking.html}
\end{itemize}

\subsection{Package-Versionen}\label{package-versionen}

Die auf den Servern des StBA und FDZ verfügbaren Packageversionen lassen
sich aus organisatorischen Gründen nicht flexibel ändern, da die
Genehmigung und Einrichtung neuer Packages nur auf Antrag durch unser
R-Betriebsteam erfolgt (i.d.R. innerhalb von 1-2 Tagen ). Die
Package-Versionen, die zum Zeitpunkt des in Kapitel~\ref{sec-benchmark}
dargestellten Performanz-Vergleiches verfügab waren, können Tabelle
Tabelle~\ref{tbl-pkg-versions} entnommen werden.

\begin{table*}

\caption{\label{tbl-pkg-versions}Versionen verfügbarer R-Packages}

\centering{

[!htb]

\centering\begingroup\fontsize{10}{12}\selectfont


\begin{tabular}{|>{}l|||>{}l|}
\hline
\cellcolor[HTML]{dce6f0}{\textcolor[HTML]{404040}{\textbf{Software}}} & \cellcolor[HTML]{dce6f0}{\textcolor[HTML]{404040}{\textbf{Version}}}\\
\hline
\cellcolor{white}{\textcolor{black}{SAS}} & \cellcolor[HTML]{f0f8ff}{9.4 (M8)}\\
\cellcolor{white}{\textcolor{black}{Stata}} & \cellcolor[HTML]{f0f8ff}{19 MP}\\
\cellcolor{white}{\textcolor{black}{R}} & \cellcolor[HTML]{f0f8ff}{4.5}\\
\cellcolor{white}{\textcolor{black}{\{data.table\}}} & \cellcolor[HTML]{f0f8ff}{1.17.8}\\
\cellcolor{white}{\textcolor{black}{\{readr\}}} & \cellcolor[HTML]{f0f8ff}{tba}\\
\cellcolor{white}{\textcolor{black}{\{dplyr\}}} & \cellcolor[HTML]{f0f8ff}{tba}\\
\cellcolor{white}{\textcolor{black}{\{vroom\}}} & \cellcolor[HTML]{f0f8ff}{tba}\\
\cellcolor{white}{\textcolor{black}{\{arrow\}}} & \cellcolor[HTML]{f0f8ff}{20.0.0.2}\\
\cellcolor{white}{\textcolor{black}{\{duckdb\}}} & \cellcolor[HTML]{f0f8ff}{1.3.2}\\
\cellcolor{white}{\textcolor{black}{\{polars\}}} & \cellcolor[HTML]{f0f8ff}{1.0.1}\\
\hline
\end{tabular}
\endgroup{}

}

\end{table*}%

\section{Performanzvergleich}\label{sec-benchmark}

RAM nicht aussagekräftig!!!

\subsection{Versuchsaufbau}\label{versuchsaufbau}

Um zwischen den in Abschnitt Kapitel~\ref{sec-methoden} vorgestellten
Datenverarbeitungsmethoden abzuwägen, haben wir deren Performanz in
einem systematischen Vergleich der Performance auf dem SAS-Server des
Statistischen Bundesamtes und den neuen OnSite R- und Stata-Servern des
FDZ gemessen. Weiterführende Performanzmessung sowie der vollständig
reproduzierbare Code sind auf Opencode verfügbar\footnote{www.opencode.de/URL-To-Be-Announced.}.

Für die Untersuchung wurden simulierte Einzeldaten in verschiedenen
Größen (0,1\%/1\%/10\% von den insgesamt 67 Millionen Beobachtungen des
BTP) betrachtet. Sie wurden aus öffentlich verfügbaren Informationen
generiert, um die wesentlichen Strukturen des BTP für unsere Analyse
künstlich nachzubilden. Abgrenzend zu den im FDZ bereitgestellten
\emph{Datenstrukturfiles} wurden die Testdaten nicht aus echten
Mikrodaten abgeleitet. Der simulierte Datensatz ist ausschließlich für
technische Testzwecke der nachfolgenden Anwendungsfälle bestimmt:

\begin{enumerate}
\def\labelenumi{\arabic{enumi}.}
\tightlist
\item
  \textbf{Fallzahlen} (pro Statistik und Jahr),
\item
  \textbf{Regression} (Einflussfaktoren auf Verlustvorträge).
\end{enumerate}

Um eine zuverlässige Performanz-Messung zu erhalten, wurde jede
Datenverarbeitungsmethode dreifach angewandt. Der Vergleich erfolgt
anhand der mittleren Messung der \textbf{Laufzeit} der Auswertung (real
verstrichene Zeit).

Daneben ist die Erfassung des \textbf{Arbeitsspeicher} der Auswertung
entscheidend, da ein übermäßiger Arbeitsspeicherverbrauch zu
schwerwiegenden Abbrüchen führen kann. Dennoch

\textbf{Festplattenspeicher} ist ausreichend vorhanden, aber unnötige
Belegung sollte vermieden werden. Feine Abstufungen der Zielgrößen sind
für unsere Aufgaben nicht gleichermaßen relevant. Eine Methode, die die
betrachteten Auswertungen in unter 10 Sekunden mit weniger als 4 GB
Arbeitsspeicher durchführen kann, wird bereits als optimal bewertet.
Weitere Optimierungen sind in diesem Fall nicht erforderlich. Falls
bereits die gewünschte Performanz erreicht wurde, rückt die Reduktion
des Entwicklungsaufwand in den Vordergrund. Entsprechend sind Aspekte
wie einfache Anwendbarkeit und Nachvollziehbarkeit -- selbst für
unerfahrene Nutzende -- ausschlaggebend.

\subsection{Vergleichswert in SAS und
Stata}\label{vergleichswert-in-sas-und-stata}

Klassisch werden die betrachteten Auswertungsaufgaben im Statischen
Bundesamt hauptsächlich mit SAS und im FDZ mit Stata gelöst. Um die
nachfolgenden Ergebnisse einzuordnen, haben wir die Performanzmessungen
für 1\% des simulierten BTP-Daten (etwa. 800.000 Beobachtungen) in SAS
und Stata durchgeführt.

\begin{table*}

\caption{\label{tbl-baseline}Baselines in SAS und Stata (anhand 1\% des simuliertes BTP)}

\centering{

[!htb]

\centering\begingroup\fontsize{8}{10}\selectfont


\begin{tabular}{|>{}l||>{}l||>{}r||>{}r||>{}r|}
\toprule
Software & Anwendung & Arbeitsspeicher & Laufzeit & Festplattenspeicher\\
\midrule
 & \cellcolor[HTML]{f0f8ff}{Fallzahl} & \cellcolor[HTML]{f0f8ff}{65 MB} & \cellcolor[HTML]{f0f8ff}{2.4 Min.} & \cellcolor[HTML]{f0f8ff}{33 GB}\\
\cmidrule{2-5}
\multirow[t]{-2}{*}{\raggedright\arraybackslash \cellcolor{white}{\textcolor{black}{SAS}}} & \cellcolor[HTML]{f0f8ff}{Regression} & \cellcolor[HTML]{f0f8ff}{7 MB} & \cellcolor[HTML]{f0f8ff}{2.8 Min.} & \cellcolor[HTML]{f0f8ff}{33 GB}\\
\cmidrule{1-5}
 & \cellcolor[HTML]{f0f8ff}{Fallzahl} & \cellcolor[HTML]{f0f8ff}{33 GB} & \cellcolor[HTML]{f0f8ff}{55,2 Sek.} & \cellcolor[HTML]{f0f8ff}{32 GB}\\
\cmidrule{2-5}
\multirow[t]{-2}{*}{\raggedright\arraybackslash \cellcolor{white}{\textcolor{black}{Stata}}} & \cellcolor[HTML]{f0f8ff}{Regression} & \cellcolor[HTML]{f0f8ff}{33 GB} & \cellcolor[HTML]{f0f8ff}{55,8 Sek.} & \cellcolor[HTML]{f0f8ff}{32 GB}\\
\bottomrule
\end{tabular}
\endgroup{}

}

\end{table*}%

Die Ergebnisse in Tabelle~\ref{tbl-baseline} zeigen, dass der
umfangreiche Arbeitsspeicher des FDZ OnSite-Servers nicht ausreicht, um
das ganze BTP in Stata einzulesen, sodass zwingend mit
Variablenselektion gearbeitet werden muss. SAS zeigt seine Stärke anhand
der geringen Beanspruchung des Arbeitsspeichers, benötigt aber mit etwa
2,3 Minuten pro Auswertung die doppelte Laufzeit gegenüber Stata.

\subsection{Kleiner Vergleich zwischen R, SAS und
Stata}\label{kleiner-vergleich-zwischen-r-sas-und-stata}

RAM Messung: http://adv-r.had.co.nz/memory.html

Für 1\% der simulierten BTP Daten, reduzieren die
R-Datenverarbeitungsmethode alle Zielgrößen gegenüber SAS und Stata
erheblich (siehe Tabelle~\ref{tbl-benchmark-small} und
Tabelle~\ref{tbl-benchmark-regr-small}). Selbst klassische R-Methoden
sind um den Faktor 7 schneller und benötigen nur 33 \% des
Festplattenspeichers. \{polars\}, \{duckdb\} und \{arrow\} lösen die
betrachteten Auswertungen sogar

\begin{itemize}
\tightlist
\item
  in \textbf{unter 2 Sekunden} (27-fach schneller),
\item
  mit unter \textbf{300 MB Arbeitsspeicher} -- \{duckdb\} und
  \{polars\}{[}\^{}polars-ram{]} sogar mit unter \textbf{5 MB} (10-fach
  weniger)\footnote{Lediglich in der Arbeitsspeicherbeanspruchung, ist
    SAS deutlich besser als klassische R-Methoden. Gegenüber
    hoch-performanten R-Methoden benötigt SAS jedoch den 10-fachen
    Arbeitsspeicher.}\footnote{\{arrow\} benötigt mit 274 MB deutlich
    mehr Arbeitsspeicher als \{polars\} und \{duckdb\}, aber kommt
    dennoch mit minimalen Entwicklungsumgebungen aus, bspw. einem
    standardmäßigen Laptop.} --
\item
  und nur mit \textbf{5 GB Festplattenspeicher} (6-fach weniger
  gegenüber SAS und Stata).
\end{itemize}

Unter den klassischen R-Methoden bietet \{data.table\} die beste
Performanz. Dabei ist wesentlich zu beachten, dass die Einlesefunktion
\texttt{fread} direkt die relevanten Variablen auswählt. Dadurch lassen
sich unsere Auswertungen um einen Faktor 3 beschleunigen und vor allem
der Arbeitsspeicher um einen Faktor 1.000 reduzieren.

Alle hoch-performanten Methoden sind dermaßen effizient, dass die
Datenmenge von 800.000 Beobachtungen nicht mehr ausreicht, um zwischen
den besten Methoden zu differenzieren. Daher erhöhen wir im nächsten
Experiment die Datenmenge auf 8.000.000 Beobachtungen (10\% BTP).

\begin{table*}

\caption{\label{tbl-benchmark-small}Performanz Messungen für die Fallzahlberechnung anhand 1\% des BTP}

\centering{

[!htb]

\centering\begingroup\fontsize{8}{10}\selectfont


\begin{tabular}{|>{}l|||>{}l|||>{}l|||>{}r|||>{}r|||>{}r|}
\hline
\cellcolor[HTML]{dce6f0}{\textcolor[HTML]{404040}{\textbf{Package}}} & \cellcolor[HTML]{dce6f0}{\textcolor[HTML]{404040}{\textbf{Dateiformat}}} & \cellcolor[HTML]{dce6f0}{\textcolor[HTML]{404040}{\textbf{Optimierung}}} & \cellcolor[HTML]{dce6f0}{\textcolor[HTML]{404040}{\textbf{Arbeitsspeicher}}} & \cellcolor[HTML]{dce6f0}{\textcolor[HTML]{404040}{\textbf{Laufzeit}}} & \cellcolor[HTML]{dce6f0}{\textcolor[HTML]{404040}{\textbf{Festplattenspeicher}}}\\
\hline
\cellcolor{white}{\textcolor{black}{\{polars\}}} & \cellcolor[HTML]{f0f8ff}{Parquet} & \cellcolor[HTML]{f0f8ff}{-} & \cellcolor[HTML]{f0f8ff}{3,3 MB} & \cellcolor[HTML]{f0f8ff}{1,0 Sek} & \cellcolor[HTML]{f0f8ff}{4,9 GB}\\
\cellcolor{white}{\textcolor{black}{\{duckdb\}}} & \cellcolor[HTML]{f0f8ff}{Parquet} & \cellcolor[HTML]{f0f8ff}{-} & \cellcolor[HTML]{f0f8ff}{4,8 MB} & \cellcolor[HTML]{f0f8ff}{1,5 Sek} & \cellcolor[HTML]{f0f8ff}{4,9 GB}\\
\cellcolor{white}{\textcolor{black}{\{arrow\}}} & \cellcolor[HTML]{f0f8ff}{Parquet} & \cellcolor[HTML]{f0f8ff}{Variablenauswahl} & \cellcolor[HTML]{f0f8ff}{273,7 MB} & \cellcolor[HTML]{f0f8ff}{1,7 Sek} & \cellcolor[HTML]{f0f8ff}{4,9 GB}\\
\cellcolor{white}{\textcolor{black}{\{data.table\}}} & \cellcolor[HTML]{f0f8ff}{CSV} & \cellcolor[HTML]{f0f8ff}{Variablenauswahl} & \cellcolor[HTML]{f0f8ff}{42,9 MB} & \cellcolor[HTML]{f0f8ff}{2,5 Sek} & \cellcolor[HTML]{f0f8ff}{13,0 GB}\\
\cellcolor{white}{\textcolor{black}{\{arrow\}}} & \cellcolor[HTML]{f0f8ff}{Parquet} & \cellcolor[HTML]{f0f8ff}{Datenaufteilung (gleichmäßig)} & \cellcolor[HTML]{f0f8ff}{229,7 MB} & \cellcolor[HTML]{f0f8ff}{2,8 Sek} & \cellcolor[HTML]{f0f8ff}{4,9 GB}\\
\cellcolor{white}{\textcolor{black}{\{arrow\}}} & \cellcolor[HTML]{f0f8ff}{Parquet} & \cellcolor[HTML]{f0f8ff}{-} & \cellcolor[HTML]{f0f8ff}{272,8 MB} & \cellcolor[HTML]{f0f8ff}{2,8 Sek} & \cellcolor[HTML]{f0f8ff}{4,9 GB}\\
\cellcolor{white}{\textcolor{black}{\{arrow\}}} & \cellcolor[HTML]{f0f8ff}{Parquet} & \cellcolor[HTML]{f0f8ff}{-} & \cellcolor[HTML]{f0f8ff}{230,4 MB} & \cellcolor[HTML]{f0f8ff}{3,8 Sek} & \cellcolor[HTML]{f0f8ff}{4,9 GB}\\
\cellcolor{white}{\textcolor{black}{\{arrow\}}} & \cellcolor[HTML]{f0f8ff}{Parquet} & \cellcolor[HTML]{f0f8ff}{Datenaufteilung (Berichtsjahr)} & \cellcolor[HTML]{f0f8ff}{229,9 MB} & \cellcolor[HTML]{f0f8ff}{3,9 Sek} & \cellcolor[HTML]{f0f8ff}{4,8 GB}\\
\cellcolor{white}{\textcolor{black}{\{data.table\}}} & \cellcolor[HTML]{f0f8ff}{CSV} & \cellcolor[HTML]{f0f8ff}{-} & \cellcolor[HTML]{f0f8ff}{33,7 GB} & \cellcolor[HTML]{f0f8ff}{7,4 Sek} & \cellcolor[HTML]{f0f8ff}{13,0 GB}\\
\cellcolor{white}{\textcolor{black}{\{data.table\}}} & \cellcolor[HTML]{f0f8ff}{CSV} & \cellcolor[HTML]{f0f8ff}{Datenaufteilung (Berichtsjahr)} & \cellcolor[HTML]{f0f8ff}{50,4 GB} & \cellcolor[HTML]{f0f8ff}{13,6 Sek} & \cellcolor[HTML]{f0f8ff}{10,5 GB}\\
\hline
\end{tabular}
\endgroup{}

}

\end{table*}%

\begin{table*}

\caption{\label{tbl-benchmark-regr-small}Performanz Messungen für die Regressionsberechnung anhand 1\% des BTP}

\centering{

[!htb]

\centering\begingroup\fontsize{8}{10}\selectfont


\begin{tabular}{|>{}l|||>{}l|||>{}l|||>{}r|||>{}r|||>{}r|}
\hline
\cellcolor[HTML]{dce6f0}{\textcolor[HTML]{404040}{\textbf{Package}}} & \cellcolor[HTML]{dce6f0}{\textcolor[HTML]{404040}{\textbf{Dateiformat}}} & \cellcolor[HTML]{dce6f0}{\textcolor[HTML]{404040}{\textbf{Optimierung}}} & \cellcolor[HTML]{dce6f0}{\textcolor[HTML]{404040}{\textbf{Arbeitsspeicher}}} & \cellcolor[HTML]{dce6f0}{\textcolor[HTML]{404040}{\textbf{Laufzeit}}} & \cellcolor[HTML]{dce6f0}{\textcolor[HTML]{404040}{\textbf{Festplattenspeicher}}}\\
\hline
\cellcolor{white}{\textcolor{black}{\{arrow\}}} & \cellcolor[HTML]{f0f8ff}{Parquet} & \cellcolor[HTML]{f0f8ff}{Datenaufteilung (Berichtsjahr)} & \cellcolor[HTML]{f0f8ff}{529,7 MB} & \cellcolor[HTML]{f0f8ff}{3,8 Sek} & \cellcolor[HTML]{f0f8ff}{4,8 GB}\\
\cellcolor{white}{\textcolor{black}{\{dplyr\}}} & \cellcolor[HTML]{f0f8ff}{CSV} & \cellcolor[HTML]{f0f8ff}{Variablenauswahl} & \cellcolor[HTML]{f0f8ff}{154,8 MB} & \cellcolor[HTML]{f0f8ff}{13,2 Sek} & \cellcolor[HTML]{f0f8ff}{13,0 GB}\\
\hline
\end{tabular}
\endgroup{}

}

\end{table*}%

\clearpage

\subsection{Größere Datenmengen in
R}\label{gruxf6uxdfere-datenmengen-in-r}

Auch für 10 \% der simulierten BTP-Daten, kann R die beiden betrachteten
Auswertungen in Echtzeit mit minimalem Arbeitsspeicher beantworten
(siehe Tabelle~\ref{tbl-benchmark} und
Tabelle~\ref{tbl-benchmark-regr}). Die Auswertungsmethoden mit der
höchsten Performanz nutzen alle das \{parquet\} und erreichen:

\begin{enumerate}
\def\labelenumi{\arabic{enumi}.}
\tightlist
\item
  \{arrow\} benötigt nur \textbf{4,7 Sekunden} und 230 MB
  Arbeitsspeicher,
\item
  \{duckdb\} benötigt nur 5,1 Sekunden und \textbf{4,7 MB
  Arbeitsspeicher}.
\end{enumerate}

{[}tba: Plot der Performanz unterschiede{]}

Unerwarteterweise ist die Laufzeit von \{polars\} für diese Datenmenge
mit 1,2 Minuten weit abgeschlagen, obwohl das gleiche Auswertungsskript
für kleinere Datenmengen immer die geringste Laufzeit aufwies. Da
\{polars\} das jüngste der betrachteten R-Packages ist, vermuten wir
Performanzprobleme in der vorliegenden Version 1.0.1\footnote{Das
  R-Package \{polars\} Version 1.0.1 war nur der erste Minor Release
  nach einer kompletten Überarbeitung. Inzwischen sind 5 neuere Major
  Releases mit diversen Performanz-Verbesserungen verfügbar.}. Zudem hat
es noch größere Entwicklungsaufwände beansprucht, da der an Python
angelehnte Syntaxstil deutlich von üblicher Programmierung in R
abweicht. Öffentliche Hilfestellung bezieht sich hauptsächlich auf
Python (oder Rust), was sich in den Antworten der KI-Chatassistenten
wiederspiegelt.

Klassische Methoden kommen bei der betrachteten Datenmenge bereits an
ihre Grenzen. Falls in \{data.table\} unbedacht der vollständige
Datensatz eingelesen wird, werden in unserem Anwemdungsfall die
verfügbaren 512 GB Arbeitsspeicher vollständig beansprucht. Schon wenn
klassische Methoden auf Datenmengen von ca. 25-33\% des verfügbaren
Arbeitsspeichers angewandt werden, ist somit deutlich mehr Erfahrung der
Nutzenden notwendig, um die relevanten Variablen manuell geschickt zu
selektieren unter Beobachtung des Arbeitsspeicherverbrauchs.

Somit kommen wir insgesamt zu dem Ergebnis, dass \{arrow\} kombiniert
mit \{parquet\} für die Bewältigung unserer Aufgaben die beste
Datenverarbeitungsmethode ist, da die Methodik hervorragende Performanz
und zugleich hohe Nutzerfreundlichkeit bietet. Im Detail haben uns der
folgenden Mehrwerte überzeugt:

\begin{enumerate}
\def\labelenumi{\arabic{enumi}.}
\tightlist
\item
  Für die relevanteste Datenmenge von 8 Mio. Beobachtungen zeigte
  \{arrow\} die geringsten Laufzeiten unter 5 Sekunden, was der
  Fachstatistik Antworten auf Sonderanfragen in Echtzeit ermöglicht und
  der Wissenschaft alle denkenswerten explorativen Analysen auszuweiten.
\item
  auch die Nutzererfahrung gefiel uns mit \{arrow\} am besten und somit
  waren die einhergehenden Entwicklungaufwände gegenüber den Anderen
  Methoden am geringsten.
\item
  ,\{arrow\} verwendet die anschauliche \{tidyverse\} Syntax nativ.
  Vorhandener Code lässt sich zum Großteil änderungsfrei und ohne
  Performanzverluste auf \{arrow\} übertragen. \{duckdb\} kann bei
  Kenntnissen in SQL zu bevorzugen sein und \{polars\} mit Erfahrungen
  in Python.
\item
  Beide Anwendungen ließen sich in unter 20 nachvollziehbaren Zeilen in
  \{arrow\} programmieren.\footnote{www.opencode.de/URL-To-Be-Announced.}
\item
  Obwohl nicht alle Features aus \{tidyverse\} für \{arrow\} verfügbar
  sind, ließ sich bisher immer eine Lösung finden, um die gewünschte
  Auswertung umzusetzen, bspw. durch Kombination mit \{duckdb\} oder
  andere Funktionsaufrufe.
\item
  \{arrow\} spart ausreichend Arbeitsspeicher (mit Verbrauch unter 1 GB)
  (tba: Fußnote zu duckdb/polars). Der Arbeitsspeicherverbrauch von
  \{arrow\} ist in unserer Auswertung nahezu konstant abhängig von der
  Datenmenge, sodass \{arrow\} leicht auf noch größere Datenmengen
  skalieren kann
\item
  Dadurch das \{arrow\} die Nutzung auch von mehreren \{parquet\}
  Dateien unterstützt spart gegenüber CSV mind. 50\% der Festplatte (für
  das reale BTP sogar 80\%+) und ermöglicht schnelle Nutzung auch von
  Datenmengen überhalb des verfügbaren Arbeitsspeichers.
\end{enumerate}

\begin{table*}

\caption{\label{tbl-benchmark}Performanz Messungen für die Fallzahlberechnung anhand 10\% des BTP}

\centering{

[!htb]

\centering\begingroup\fontsize{8}{10}\selectfont


\begin{tabular}{|>{}l|||>{}l|||>{}l|||>{}r|||>{}r|||>{}r|}
\hline
\cellcolor[HTML]{dce6f0}{\textcolor[HTML]{404040}{\textbf{Package}}} & \cellcolor[HTML]{dce6f0}{\textcolor[HTML]{404040}{\textbf{Dateiformat}}} & \cellcolor[HTML]{dce6f0}{\textcolor[HTML]{404040}{\textbf{Optimierung}}} & \cellcolor[HTML]{dce6f0}{\textcolor[HTML]{404040}{\textbf{Arbeitsspeicher}}} & \cellcolor[HTML]{dce6f0}{\textcolor[HTML]{404040}{\textbf{Laufzeit}}} & \cellcolor[HTML]{dce6f0}{\textcolor[HTML]{404040}{\textbf{Festplattenspeicher}}}\\
\hline
\cellcolor{white}{\textcolor{black}{\{arrow\}}} & \cellcolor[HTML]{f0f8ff}{Parquet} & \cellcolor[HTML]{f0f8ff}{Datenaufteilung (gleichmäßig)} & \cellcolor[HTML]{f0f8ff}{229,9 MB} & \cellcolor[HTML]{f0f8ff}{4,7 Sek} & \cellcolor[HTML]{f0f8ff}{48,8 GB}\\
\cellcolor{white}{\textcolor{black}{\{duckdb\}}} & \cellcolor[HTML]{f0f8ff}{Parquet} & \cellcolor[HTML]{f0f8ff}{-} & \cellcolor[HTML]{f0f8ff}{4,7 MB} & \cellcolor[HTML]{f0f8ff}{5,1 Sek} & \cellcolor[HTML]{f0f8ff}{49,4 GB}\\
\cellcolor{white}{\textcolor{black}{\{arrow\}}} & \cellcolor[HTML]{f0f8ff}{Parquet} & \cellcolor[HTML]{f0f8ff}{Variablenauswahl} & \cellcolor[HTML]{f0f8ff}{702,1 MB} & \cellcolor[HTML]{f0f8ff}{7,5 Sek} & \cellcolor[HTML]{f0f8ff}{49,4 GB}\\
\cellcolor{white}{\textcolor{black}{\{arrow\}}} & \cellcolor[HTML]{f0f8ff}{Parquet} & \cellcolor[HTML]{f0f8ff}{Datenaufteilung (Berichtsjahr)} & \cellcolor[HTML]{f0f8ff}{4,7 MB} & \cellcolor[HTML]{f0f8ff}{11,7 Sek} & \cellcolor[HTML]{f0f8ff}{48,1 GB}\\
\cellcolor{white}{\textcolor{black}{\{arrow\}}} & \cellcolor[HTML]{f0f8ff}{Parquet} & \cellcolor[HTML]{f0f8ff}{-} & \cellcolor[HTML]{f0f8ff}{230,4 MB} & \cellcolor[HTML]{f0f8ff}{17,0 Sek} & \cellcolor[HTML]{f0f8ff}{49,4 GB}\\
\cellcolor{white}{\textcolor{black}{\{arrow\}}} & \cellcolor[HTML]{f0f8ff}{Parquet} & \cellcolor[HTML]{f0f8ff}{-} & \cellcolor[HTML]{f0f8ff}{701,2 MB} & \cellcolor[HTML]{f0f8ff}{20,6 Sek} & \cellcolor[HTML]{f0f8ff}{49,4 GB}\\
\cellcolor{white}{\textcolor{black}{\{data.table\}}} & \cellcolor[HTML]{f0f8ff}{CSV} & \cellcolor[HTML]{f0f8ff}{Variablenauswahl} & \cellcolor[HTML]{f0f8ff}{427,5 MB} & \cellcolor[HTML]{f0f8ff}{23,6 Sek} & \cellcolor[HTML]{f0f8ff}{130,2 GB}\\
\cellcolor{white}{\textcolor{black}{\{polars\}}} & \cellcolor[HTML]{f0f8ff}{Parquet} & \cellcolor[HTML]{f0f8ff}{-} & \cellcolor[HTML]{f0f8ff}{3,3 MB} & \cellcolor[HTML]{f0f8ff}{1,2 Min} & \cellcolor[HTML]{f0f8ff}{49,4 GB}\\
\cellcolor{white}{\textcolor{black}{\{data.table\}}} & \cellcolor[HTML]{f0f8ff}{CSV} & \cellcolor[HTML]{f0f8ff}{-} & \cellcolor[HTML]{f0f8ff}{336,0 GB} & \cellcolor[HTML]{f0f8ff}{2,2 Min} & \cellcolor[HTML]{f0f8ff}{130,2 GB}\\
\cellcolor{white}{\textcolor{black}{\{data.table\}}} & \cellcolor[HTML]{f0f8ff}{CSV} & \cellcolor[HTML]{f0f8ff}{Datenaufteilung (Berichtsjahr)} & \cellcolor[HTML]{f0f8ff}{506,4 GB} & \cellcolor[HTML]{f0f8ff}{5,0 Min} & \cellcolor[HTML]{f0f8ff}{104,6 GB}\\
\hline
\end{tabular}
\endgroup{}

}

\end{table*}%

\begin{table*}

\caption{\label{tbl-benchmark-regr}Performanz Messungen für die Regressionsberechnung anhand 10\% des BTP}

\centering{

[!htb]

\centering\begingroup\fontsize{8}{10}\selectfont


\begin{tabular}{|>{}l|||>{}l|||>{}l|||>{}r|||>{}r|||>{}r|}
\hline
\cellcolor[HTML]{dce6f0}{\textcolor[HTML]{404040}{\textbf{Package}}} & \cellcolor[HTML]{dce6f0}{\textcolor[HTML]{404040}{\textbf{Dateiformat}}} & \cellcolor[HTML]{dce6f0}{\textcolor[HTML]{404040}{\textbf{Optimierung}}} & \cellcolor[HTML]{dce6f0}{\textcolor[HTML]{404040}{\textbf{Arbeitsspeicher}}} & \cellcolor[HTML]{dce6f0}{\textcolor[HTML]{404040}{\textbf{Laufzeit}}} & \cellcolor[HTML]{dce6f0}{\textcolor[HTML]{404040}{\textbf{Festplattenspeicher}}}\\
\hline
\cellcolor{white}{\textcolor{black}{\{arrow\}}} & \cellcolor[HTML]{f0f8ff}{Parquet} & \cellcolor[HTML]{f0f8ff}{Datenaufteilung (Berichtsjahr)} & \cellcolor[HTML]{f0f8ff}{782,1 MB} & \cellcolor[HTML]{f0f8ff}{5,2 Sek} & \cellcolor[HTML]{f0f8ff}{48,1 GB}\\
\cellcolor{white}{\textcolor{black}{\{dplyr\}}} & \cellcolor[HTML]{f0f8ff}{CSV} & \cellcolor[HTML]{f0f8ff}{Variablenauswahl} & \cellcolor[HTML]{f0f8ff}{1,5 GB} & \cellcolor[HTML]{f0f8ff}{1,9 Min} & \cellcolor[HTML]{f0f8ff}{130,2 GB}\\
\hline
\end{tabular}
\endgroup{}

}

\end{table*}%

{[}tba: Plot RAM-Bedarf abhängig von der Datenmenge{]}

\setcounter{subsection}{4}

\section{Ergebnis für das vollständige BTP}\label{sec-ergebnis}

Anwendung der besten Methode auf die echten Daten des vollständige BTP
ergibt:

\begin{itemize}
\tightlist
\item
  Performanz von \{arrow\} Fallzahlen (ein paar Sekunden) und Regression
  (Sekunden)
\item
  leicht lesbar in \{arrow\}
\end{itemize}

\subsection{Ergebnis}\label{ergebnis}

{[}Ergebnistabelle Fallzahl, ggf. als Heatmap{]}

-\textgreater{} Zukünftige Nutzung für MDR wird umgesetzt.

{[}Model Summary?{]}

\subsection{Vergleich zu SAS, (Stata),
Datatable}\label{vergleich-zu-sas-stata-datatable}

\begin{itemize}
\tightlist
\item
  Ergebnis in SAS
\item
  (optional) Datatable ohne sekeltives Einlesen führt zu RAM Crash
  (RAM-Fehler?)
\item
  (optional) Datatable Sel.
\end{itemize}

Real world application unserer Ergebnisse auf das Echte BTP

Ergebnis für gesamten Datensatz

-\textgreater{} Größere Speicherreduktion, da das echte BTP mehr
fehlende Werte enthält, als anhand der externen Daten angenommen.

\setcounter{subsection}{5}

\section{Fazit}\label{sec-ausblick}

\subsection{Kernaussage des Artikels}\label{kernaussage-des-artikels}

Wir empfehlen eine Datenhaltung im Parquet-Format und \textbf{R}
incl.~\{arrow\} + \{parquet\} (bei Bedarf \{duckdb\}, \{polars\}) zu
erproben und zu etablieren! Denn unsere Untersuchung demonstriert ganz
neue Möglichkeiten, analytische Fragestellungen anhand großer
Datenmengen in Echtzeit zu beantworten.

\subsection{Auswirkung auf IT-Infrastruktur und
Beratung}\label{auswirkung-auf-it-infrastruktur-und-beratung}

\begin{itemize}
\tightlist
\item
  Fokus auf Hardware und Software mit größtem Mehrwert (weniger RAM und
  R anstatt Spark, GPUs, Datenbanken, Cloud, \ldots)
\item
  Fokus auf Maintainance und Patching von R, Parquet und hier genannte
  Packages \{arrow\}, \{duckdb\}, \{polars\}
\item
  Etablierung dieser Softwarelösungen kann langfristig\ldots{}

  \begin{itemize}
  \tightlist
  \item
    die Hardwareanforderungen reduzieren
  \item
    die Analysesoftwarevielfalt verschlanken, somit Beschaffungen
    einsparen (besserer Ausdruck?)
  \end{itemize}
\item
  Schulungsangebot benötigt für \emph{``hoch-performante Auswertungen in
  R''}
\item
  eigener Beitrag zu Open Source entscheidend, um die Etablierung zu
  fördern.
\end{itemize}

\subsection{Auswirkungen auf die
Statistikproduktion}\label{auswirkungen-auf-die-statistikproduktion}

\begin{itemize}
\tightlist
\item
  Steigerung der Effizienz und Reaktionsgeschwindigkeit
\item
  (Vorsichtig!? Annette, lieber weg oder rein?) Vereinfachung der
  Statistikproduktion, durch leicht nachvollziehbaren Code und Fokus auf
  Wissensaufbau in R
\item
  Aufbau eines Datenbestands in Parquet
\item
  Anwendung auf weitere Produkte

  \begin{itemize}
  \tightlist
  \item
    Überarbeitung des BTP
  \item
    Mindeststeuer (oder?)
  \item
    DRG
  \item
    TPP 2
  \end{itemize}
\end{itemize}

\subsection{Auswirkungen auf die FDZ}\label{auswirkungen-auf-die-fdz}

\begin{itemize}
\tightlist
\item
  Berücksichtigung in Mustersyntaxen und Nutzendenberatung im FDZ Bund.
  Teilen von Code-Snippets und Templates fördern.
\item
  mögliche Vereinfachung der Betreuung (Fokus auf R, Reduktion von
  Performanz- und Programmierproblemen)
\item
  Bereitstellung von Parquet in geeigneter Aufteilung empfohlen
\item
  Erprobung im FDZ Bund Vorbild für andere FDZ
\item
  Optimierung von Hardware und Software Angebot
\end{itemize}

\subsection{Auswirkung auf die
Wissenschaft}\label{auswirkung-auf-die-wissenschaft}

\begin{itemize}
\tightlist
\item
  Verwendung von R bietet am FDZ Bund einen deutlichen
  \textbf{Wettbewerbsvorteil in der emp. Forschung}. Lediglich R ist (am
  GWAP und via KDFV) am FDZ Bund in der Lage große Datenmengen innerhalb
  von Sekunden zu analysieren und Datenmengen überhalb des verfügbaren
  Arbeitspeicher, wie das BTP, ohne vorherige Variablenselektion
  auszuwerten.

  \begin{itemize}
  \tightlist
  \item
    Vollmaterial kann statt Stichproben angeboten werden. Erleichtert
    die FDZ Betreuung und GWAP/KDFV-Verfügbarkeit
  \item
    mehr explorative oder komplexe Analysen sind in kürzerer Zeit
    möglich.
  \end{itemize}
\item
  Dazu ist R als frei verfügbare Open-Source Software leicht zu lernen
  und \{arrow\} bietet hoch-performante Analysen selbst für
  R-Basiskenntnisse. Empfehlung an die emp. Wissenschaft, sich verstärkt
  mit R zu befassen.
\item
  Für analoge über-RAM-Verarbeitungen bietet Stata aktuell keine Lösung.
  SAS bietet bei weitem nicht hier gemessene Performanz von R. R
  hervorragende Möglichkeiten bietet Performanz und
  Anwenderfreundlichkeit zu übertreffen.
\item
  Weiteres Vorgehen: \textbf{NeSt-Pilotprojekt}. Aus dem realen Use Case
  Optimierung der Systemkonfiguration und Beratungsbedarfe
  identifizieren.
\end{itemize}

\subsection{Schlussbemerkung}\label{schlussbemerkung}

Die Analyse großer Datensätze in der amtlichen Statstik und den
Forschungsdatenzen können von einem signifikanten Performanz-Boost
profitieren, wenn die vorgestellten Werkzeuge adaptiert werden.
\textbf{Arrow, DuckDB, Polars und Parquet} -- bieten das Potenzial, die
Effizienz erheblich zu steigern und neue Analysemöglichkeiten zu
eröffnen.

Der Übergang erfordert jedoch eine \textbf{koordinierte Anstrengung} von
Infrastrukturbetreibern, amtlicher Statistik, Forschungsdatenzentren und
Forschenden. Die systematische Wissensaufbau, Dokumentation und
Verbreitung von Best Practices wird entscheidend für den Erfolg sein.

Die amtliche Statistik kann von diesen Entwicklungen in mehrfacher
Hinsicht profitieren: \textbf{schnellere Erkenntnisgewinnung},
\textbf{umfassendere Analysen} und \textbf{effizientere
Ressourcennutzung}. Dies trägt letztlich zur Stärkung der
evidenzbasierten Politik- und Gesellschaftsberatung bei.

\section*{Anhang}\label{anhang}
\addcontentsline{toc}{section}{Anhang}

Inhalte für die Veröffentlichung in Open Code.

\subsection*{Weitere Performancemessungen für kleinere
Datenmengen.}\label{weitere-performancemessungen-fuxfcr-kleinere-datenmengen.}
\addcontentsline{toc}{subsection}{Weitere Performancemessungen für
kleinere Datenmengen.}

\begin{table*}

\caption{\label{tbl-bm-cnt-small}Benchmarkergebnisse für Fallzahlen (kleine Stichproben des BTP)}

\centering{

[!htb]

\centering\begingroup\fontsize{8}{10}\selectfont


\begin{tabular}{|>{}r||>{}l||>{}l||>{}l||>{}r||>{}r||>{}r|}
\toprule
\cellcolor[HTML]{dce6f0}{\textcolor[HTML]{404040}{\textbf{Beob.}}} & \cellcolor[HTML]{dce6f0}{\textcolor[HTML]{404040}{\textbf{Package}}} & \cellcolor[HTML]{dce6f0}{\textcolor[HTML]{404040}{\textbf{Dateiformat}}} & \cellcolor[HTML]{dce6f0}{\textcolor[HTML]{404040}{\textbf{Optimierung}}} & \cellcolor[HTML]{dce6f0}{\textcolor[HTML]{404040}{\textbf{Arbeitsspeicher}}} & \cellcolor[HTML]{dce6f0}{\textcolor[HTML]{404040}{\textbf{Laufzeit}}} & \cellcolor[HTML]{dce6f0}{\textcolor[HTML]{404040}{\textbf{Festplattenspeicher}}}\\
\midrule
\cellcolor{white}{\textcolor{black}{10.000}} & \cellcolor[HTML]{f0f8ff}{\{polars\}} & \cellcolor[HTML]{f0f8ff}{Parquet} & \cellcolor[HTML]{f0f8ff}{-} & \cellcolor[HTML]{f0f8ff}{3,3 MB} & \cellcolor[HTML]{f0f8ff}{0,4 Sek} & \cellcolor[HTML]{f0f8ff}{492,6 MB}\\
\cellcolor{white}{\textcolor{black}{10.000}} & \cellcolor[HTML]{f0f8ff}{\{data.table\}} & \cellcolor[HTML]{f0f8ff}{CSV} & \cellcolor[HTML]{f0f8ff}{Variablenauswahl} & \cellcolor[HTML]{f0f8ff}{4,4 MB} & \cellcolor[HTML]{f0f8ff}{0,7 Sek} & \cellcolor[HTML]{f0f8ff}{1,3 GB}\\
\cellcolor{white}{\textcolor{black}{10.000}} & \cellcolor[HTML]{f0f8ff}{\{arrow\}} & \cellcolor[HTML]{f0f8ff}{Parquet} & \cellcolor[HTML]{f0f8ff}{Variablenauswahl} & \cellcolor[HTML]{f0f8ff}{230,7 MB} & \cellcolor[HTML]{f0f8ff}{1,1 Sek} & \cellcolor[HTML]{f0f8ff}{492,6 MB}\\
\cellcolor{white}{\textcolor{black}{10.000}} & \cellcolor[HTML]{f0f8ff}{\{readr\}} & \cellcolor[HTML]{f0f8ff}{CSV} & \cellcolor[HTML]{f0f8ff}{Variablenauswahl} & \cellcolor[HTML]{f0f8ff}{9,3 MB} & \cellcolor[HTML]{f0f8ff}{1,2 Sek} & \cellcolor[HTML]{f0f8ff}{1,3 GB}\\
\cellcolor{white}{\textcolor{black}{10.000}} & \cellcolor[HTML]{f0f8ff}{\{arrow\}} & \cellcolor[HTML]{f0f8ff}{Parquet} & \cellcolor[HTML]{f0f8ff}{-} & \cellcolor[HTML]{f0f8ff}{229,8 MB} & \cellcolor[HTML]{f0f8ff}{1,6 Sek} & \cellcolor[HTML]{f0f8ff}{492,6 MB}\\
\cellcolor{white}{\textcolor{black}{10.000}} & \cellcolor[HTML]{f0f8ff}{\{duckdb\}} & \cellcolor[HTML]{f0f8ff}{Parquet} & \cellcolor[HTML]{f0f8ff}{-} & \cellcolor[HTML]{f0f8ff}{4,7 MB} & \cellcolor[HTML]{f0f8ff}{2,9 Sek} & \cellcolor[HTML]{f0f8ff}{492,6 MB}\\
\cellcolor{white}{\textcolor{black}{10.000}} & \cellcolor[HTML]{f0f8ff}{\{data.table\}} & \cellcolor[HTML]{f0f8ff}{CSV} & \cellcolor[HTML]{f0f8ff}{-} & \cellcolor[HTML]{f0f8ff}{3,4 GB} & \cellcolor[HTML]{f0f8ff}{3,6 Sek} & \cellcolor[HTML]{f0f8ff}{1,3 GB}\\
\cellcolor{white}{\textcolor{black}{10.000}} & \cellcolor[HTML]{f0f8ff}{\{arrow\}} & \cellcolor[HTML]{f0f8ff}{Parquet} & \cellcolor[HTML]{f0f8ff}{Datenaufteilung (Berichtsjahr)} & \cellcolor[HTML]{f0f8ff}{229,9 MB} & \cellcolor[HTML]{f0f8ff}{4,2 Sek} & \cellcolor[HTML]{f0f8ff}{486,9 MB}\\
\cellcolor{white}{\textcolor{black}{10.000}} & \cellcolor[HTML]{f0f8ff}{\{arrow\}} & \cellcolor[HTML]{f0f8ff}{Parquet} & \cellcolor[HTML]{f0f8ff}{-} & \cellcolor[HTML]{f0f8ff}{230,4 MB} & \cellcolor[HTML]{f0f8ff}{4,6 Sek} & \cellcolor[HTML]{f0f8ff}{492,6 MB}\\
\cellcolor{white}{\textcolor{black}{10.000}} & \cellcolor[HTML]{f0f8ff}{\{data.table\}} & \cellcolor[HTML]{f0f8ff}{CSV} & \cellcolor[HTML]{f0f8ff}{Datenaufteilung (Berichtsjahr)} & \cellcolor[HTML]{f0f8ff}{5,1 GB} & \cellcolor[HTML]{f0f8ff}{4,8 Sek} & \cellcolor[HTML]{f0f8ff}{1,0 GB}\\
\cellcolor{white}{\textcolor{black}{10.000}} & \cellcolor[HTML]{f0f8ff}{\{arrow\}} & \cellcolor[HTML]{f0f8ff}{Parquet} & \cellcolor[HTML]{f0f8ff}{Datenaufteilung (gleichmäßig)} & \cellcolor[HTML]{f0f8ff}{229,9 MB} & \cellcolor[HTML]{f0f8ff}{7,1 Sek} & \cellcolor[HTML]{f0f8ff}{483,5 MB}\\
\cellcolor{white}{\textcolor{black}{10.000}} & \cellcolor[HTML]{f0f8ff}{\{readr\}} & \cellcolor[HTML]{f0f8ff}{CSV} & \cellcolor[HTML]{f0f8ff}{-} & \cellcolor[HTML]{f0f8ff}{1,7 GB} & \cellcolor[HTML]{f0f8ff}{8,6 Sek} & \cellcolor[HTML]{f0f8ff}{1,3 GB}\\
\cellcolor{white}{\textcolor{black}{10.000}} & \cellcolor[HTML]{f0f8ff}{\{arrow\}} & \cellcolor[HTML]{f0f8ff}{CSV} & \cellcolor[HTML]{f0f8ff}{-} & \cellcolor[HTML]{f0f8ff}{27,4 MB} & \cellcolor[HTML]{f0f8ff}{15,7 Sek} & \cellcolor[HTML]{f0f8ff}{1,3 GB}\\
\cellcolor{white}{\textcolor{black}{10.000}} & \cellcolor[HTML]{f0f8ff}{\{data.table\}} & \cellcolor[HTML]{f0f8ff}{CSV.GZ} & \cellcolor[HTML]{f0f8ff}{-} & \cellcolor[HTML]{f0f8ff}{4,7 GB} & \cellcolor[HTML]{f0f8ff}{21,2 Sek} & \cellcolor[HTML]{f0f8ff}{421,9 MB}\\
\cellcolor{white}{\textcolor{black}{10.000}} & \cellcolor[HTML]{f0f8ff}{\{base\}} & \cellcolor[HTML]{f0f8ff}{CSV} & \cellcolor[HTML]{f0f8ff}{-} & \cellcolor[HTML]{f0f8ff}{9,7 GB} & \cellcolor[HTML]{f0f8ff}{4,9 Min} & \cellcolor[HTML]{f0f8ff}{1,3 GB}\\
\cellcolor{white}{\textcolor{black}{1.000}} & \cellcolor[HTML]{f0f8ff}{\{data.table\}} & \cellcolor[HTML]{f0f8ff}{CSV} & \cellcolor[HTML]{f0f8ff}{Variablenauswahl} & \cellcolor[HTML]{f0f8ff}{0,6 MB} & \cellcolor[HTML]{f0f8ff}{0,1 Sek} & \cellcolor[HTML]{f0f8ff}{129,8 MB}\\
\cellcolor{white}{\textcolor{black}{1.000}} & \cellcolor[HTML]{f0f8ff}{\{readr\}} & \cellcolor[HTML]{f0f8ff}{CSV} & \cellcolor[HTML]{f0f8ff}{Variablenauswahl} & \cellcolor[HTML]{f0f8ff}{2,7 MB} & \cellcolor[HTML]{f0f8ff}{0,2 Sek} & \cellcolor[HTML]{f0f8ff}{129,8 MB}\\
\cellcolor{white}{\textcolor{black}{1.000}} & \cellcolor[HTML]{f0f8ff}{\{polars\}} & \cellcolor[HTML]{f0f8ff}{Parquet} & \cellcolor[HTML]{f0f8ff}{-} & \cellcolor[HTML]{f0f8ff}{3,3 MB} & \cellcolor[HTML]{f0f8ff}{0,4 Sek} & \cellcolor[HTML]{f0f8ff}{48,1 MB}\\
\cellcolor{white}{\textcolor{black}{1.000}} & \cellcolor[HTML]{f0f8ff}{\{data.table\}} & \cellcolor[HTML]{f0f8ff}{CSV} & \cellcolor[HTML]{f0f8ff}{-} & \cellcolor[HTML]{f0f8ff}{341,8 MB} & \cellcolor[HTML]{f0f8ff}{0,9 Sek} & \cellcolor[HTML]{f0f8ff}{129,8 MB}\\
\cellcolor{white}{\textcolor{black}{1.000}} & \cellcolor[HTML]{f0f8ff}{\{duckdb\}} & \cellcolor[HTML]{f0f8ff}{Parquet} & \cellcolor[HTML]{f0f8ff}{-} & \cellcolor[HTML]{f0f8ff}{4,7 MB} & \cellcolor[HTML]{f0f8ff}{1,1 Sek} & \cellcolor[HTML]{f0f8ff}{48,1 MB}\\
\cellcolor{white}{\textcolor{black}{1.000}} & \cellcolor[HTML]{f0f8ff}{\{arrow\}} & \cellcolor[HTML]{f0f8ff}{Parquet} & \cellcolor[HTML]{f0f8ff}{Variablenauswahl} & \cellcolor[HTML]{f0f8ff}{226,2 MB} & \cellcolor[HTML]{f0f8ff}{1,2 Sek} & \cellcolor[HTML]{f0f8ff}{48,1 MB}\\
\cellcolor{white}{\textcolor{black}{1.000}} & \cellcolor[HTML]{f0f8ff}{\{arrow\}} & \cellcolor[HTML]{f0f8ff}{Parquet} & \cellcolor[HTML]{f0f8ff}{-} & \cellcolor[HTML]{f0f8ff}{225,3 MB} & \cellcolor[HTML]{f0f8ff}{1,4 Sek} & \cellcolor[HTML]{f0f8ff}{48,1 MB}\\
\cellcolor{white}{\textcolor{black}{1.000}} & \cellcolor[HTML]{f0f8ff}{\{arrow\}} & \cellcolor[HTML]{f0f8ff}{CSV} & \cellcolor[HTML]{f0f8ff}{-} & \cellcolor[HTML]{f0f8ff}{3,5 MB} & \cellcolor[HTML]{f0f8ff}{1,5 Sek} & \cellcolor[HTML]{f0f8ff}{129,8 MB}\\
\cellcolor{white}{\textcolor{black}{1.000}} & \cellcolor[HTML]{f0f8ff}{\{data.table\}} & \cellcolor[HTML]{f0f8ff}{CSV.GZ} & \cellcolor[HTML]{f0f8ff}{-} & \cellcolor[HTML]{f0f8ff}{491,7 MB} & \cellcolor[HTML]{f0f8ff}{2,7 Sek} & \cellcolor[HTML]{f0f8ff}{42,1 MB}\\
\cellcolor{white}{\textcolor{black}{1.000}} & \cellcolor[HTML]{f0f8ff}{\{data.table\}} & \cellcolor[HTML]{f0f8ff}{CSV} & \cellcolor[HTML]{f0f8ff}{Datenaufteilung (Berichtsjahr)} & \cellcolor[HTML]{f0f8ff}{510,9 MB} & \cellcolor[HTML]{f0f8ff}{2,8 Sek} & \cellcolor[HTML]{f0f8ff}{104,5 MB}\\
\cellcolor{white}{\textcolor{black}{1.000}} & \cellcolor[HTML]{f0f8ff}{\{arrow\}} & \cellcolor[HTML]{f0f8ff}{Parquet} & \cellcolor[HTML]{f0f8ff}{Datenaufteilung (Berichtsjahr)} & \cellcolor[HTML]{f0f8ff}{229,9 MB} & \cellcolor[HTML]{f0f8ff}{2,9 Sek} & \cellcolor[HTML]{f0f8ff}{53,5 MB}\\
\cellcolor{white}{\textcolor{black}{1.000}} & \cellcolor[HTML]{f0f8ff}{\{arrow\}} & \cellcolor[HTML]{f0f8ff}{Parquet} & \cellcolor[HTML]{f0f8ff}{-} & \cellcolor[HTML]{f0f8ff}{230,4 MB} & \cellcolor[HTML]{f0f8ff}{3,7 Sek} & \cellcolor[HTML]{f0f8ff}{48,1 MB}\\
\cellcolor{white}{\textcolor{black}{1.000}} & \cellcolor[HTML]{f0f8ff}{\{arrow\}} & \cellcolor[HTML]{f0f8ff}{Parquet} & \cellcolor[HTML]{f0f8ff}{Datenaufteilung (gleichmäßig)} & \cellcolor[HTML]{f0f8ff}{229,7 MB} & \cellcolor[HTML]{f0f8ff}{4,0 Sek} & \cellcolor[HTML]{f0f8ff}{50,5 MB}\\
\cellcolor{white}{\textcolor{black}{1.000}} & \cellcolor[HTML]{f0f8ff}{\{readr\}} & \cellcolor[HTML]{f0f8ff}{CSV} & \cellcolor[HTML]{f0f8ff}{-} & \cellcolor[HTML]{f0f8ff}{197,4 MB} & \cellcolor[HTML]{f0f8ff}{4,9 Sek} & \cellcolor[HTML]{f0f8ff}{129,8 MB}\\
\cellcolor{white}{\textcolor{black}{1.000}} & \cellcolor[HTML]{f0f8ff}{\{base\}} & \cellcolor[HTML]{f0f8ff}{CSV} & \cellcolor[HTML]{f0f8ff}{-} & \cellcolor[HTML]{f0f8ff}{713,4 MB} & \cellcolor[HTML]{f0f8ff}{23,3 Sek} & \cellcolor[HTML]{f0f8ff}{129,8 MB}\\
\bottomrule
\end{tabular}
\endgroup{}

}

\end{table*}%

\begin{table*}

\caption{\label{tbl-bm-reg-small}Benchmarkergebnisse für Regression (kleine Stichproben des BTP)}

\centering{

[!htb]

\centering\begingroup\fontsize{8}{10}\selectfont


\begin{tabular}{|>{}r||>{}l||>{}l||>{}l||>{}r||>{}r||>{}r|}
\toprule
\cellcolor[HTML]{dce6f0}{\textcolor[HTML]{404040}{\textbf{Beob.}}} & \cellcolor[HTML]{dce6f0}{\textcolor[HTML]{404040}{\textbf{Package}}} & \cellcolor[HTML]{dce6f0}{\textcolor[HTML]{404040}{\textbf{Dateiformat}}} & \cellcolor[HTML]{dce6f0}{\textcolor[HTML]{404040}{\textbf{Optimierung}}} & \cellcolor[HTML]{dce6f0}{\textcolor[HTML]{404040}{\textbf{Arbeitsspeicher}}} & \cellcolor[HTML]{dce6f0}{\textcolor[HTML]{404040}{\textbf{Laufzeit}}} & \cellcolor[HTML]{dce6f0}{\textcolor[HTML]{404040}{\textbf{Festplattenspeicher}}}\\
\midrule
\cellcolor{white}{\textcolor{black}{10.000}} & \cellcolor[HTML]{f0f8ff}{\{dplyr\}} & \cellcolor[HTML]{f0f8ff}{CSV} & \cellcolor[HTML]{f0f8ff}{Variablenauswahl} & \cellcolor[HTML]{f0f8ff}{17,6 MB} & \cellcolor[HTML]{f0f8ff}{1,3 Sek} & \cellcolor[HTML]{f0f8ff}{1,3 GB}\\
\cellcolor{white}{\textcolor{black}{10.000}} & \cellcolor[HTML]{f0f8ff}{\{arrow\}} & \cellcolor[HTML]{f0f8ff}{Parquet} & \cellcolor[HTML]{f0f8ff}{Datenaufteilung (Berichtsjahr)} & \cellcolor[HTML]{f0f8ff}{459,3 MB} & \cellcolor[HTML]{f0f8ff}{5,0 Sek} & \cellcolor[HTML]{f0f8ff}{486,9 MB}\\
\cellcolor{white}{\textcolor{black}{10.000}} & \cellcolor[HTML]{f0f8ff}{\{dplyr\}} & \cellcolor[HTML]{f0f8ff}{CSV} & \cellcolor[HTML]{f0f8ff}{-} & \cellcolor[HTML]{f0f8ff}{1,8 GB} & \cellcolor[HTML]{f0f8ff}{10,8 Sek} & \cellcolor[HTML]{f0f8ff}{1,3 GB}\\
\cellcolor{white}{\textcolor{black}{1.000}} & \cellcolor[HTML]{f0f8ff}{\{dplyr\}} & \cellcolor[HTML]{f0f8ff}{CSV} & \cellcolor[HTML]{f0f8ff}{Variablenauswahl} & \cellcolor[HTML]{f0f8ff}{4,2 MB} & \cellcolor[HTML]{f0f8ff}{0,3 Sek} & \cellcolor[HTML]{f0f8ff}{129,8 MB}\\
\cellcolor{white}{\textcolor{black}{1.000}} & \cellcolor[HTML]{f0f8ff}{\{dplyr\}} & \cellcolor[HTML]{f0f8ff}{CSV} & \cellcolor[HTML]{f0f8ff}{-} & \cellcolor[HTML]{f0f8ff}{206,9 MB} & \cellcolor[HTML]{f0f8ff}{6,4 Sek} & \cellcolor[HTML]{f0f8ff}{129,8 MB}\\
\cellcolor{white}{\textcolor{black}{1.000}} & \cellcolor[HTML]{f0f8ff}{\{arrow\}} & \cellcolor[HTML]{f0f8ff}{Parquet} & \cellcolor[HTML]{f0f8ff}{Datenaufteilung (Berichtsjahr)} & \cellcolor[HTML]{f0f8ff}{452,5 MB} & \cellcolor[HTML]{f0f8ff}{6,8 Sek} & \cellcolor[HTML]{f0f8ff}{53,5 MB}\\
\bottomrule
\end{tabular}
\endgroup{}

}

\end{table*}%


\printbibliography



\end{document}
